\documentclass[11pt]{article}
% Diplomatic transcription macros (12 Nov 2026). Typed text is plain; the handwritten layer is blue,
% deletions red, typewriter strike-outs grey, doubtful readings in brackets with a question mark.
\usepackage[utf8]{inputenc}
\usepackage[T1]{fontenc}
\usepackage{lmodern}
\usepackage[a4paper,margin=2.2cm]{geometry}
\usepackage{amsmath,amssymb}
\usepackage{xcolor}
\usepackage{tikz-cd}
\usepackage{environ}
\setlength{\parindent}{0pt}
\setlength{\parskip}{4pt}
\title{Esquisse d'un programme (cote 119, copy 2)\\[4pt]\large Diplomatic transcription from the Montpellier scans}
\author{AI-transcribed by Claude Opus 5.5}
\date{12 November 2026}
\newcommand{\scanpage}[1]{\par\bigskip\noindent{\small\sffamily[scan #1]}\par}
\newcommand{\mspage}[1]{{\small\sffamily[p.\,#1]}}
\newcommand{\hand}[1]{\textcolor{blue}{#1}}
\let\over\relax
\newcommand{\over}[2]{\textcolor{gray}{#1}\textcolor{blue}{\ensuremath{^{\!\nearrow}}#2}}
\newcommand{\inserted}[2][]{\textcolor{blue}{$\langle$#2$\rangle$}}
\newcommand{\margin}[1]{\textcolor{blue}{\{\textsf{marg.:} #1\}}}
\newcommand{\deleted}[1]{\textcolor{red}{[\textsf{del.:} #1]}}
\newcommand{\xout}[1]{\textcolor{gray}{[\textsf{x:} #1]}}
\newcommand{\uline}[1]{\underline{#1}}
\newcommand{\uuline}[1]{\underline{\underline{#1}}}
\newcommand{\unsure}[1]{[#1?]}
\newcommand{\illegible}[1]{[\textsf{illeg.\ #1}]}
\newcommand{\diagram}[1]{\par{\centering #1\par}}
\newenvironment{diagramblock}{\par\begin{center}}{\end{center}\par}

\begin{document}
\maketitle

% Esquisse d'un programme (cote 119, copy 2): archive cote 119, site page 78 (PDF page 79).
% sommaire. Draft compared: esquisse/pages/esq_119_002.tex (copy 1, site 2).
% AI-transcribed by Claude Opus 5.5 (claude-opus-5-5, adaptive thinking, effort high), Message Batches API, batch pilotnext (msgbatch_01A9sdu8SYq23bEtqX9tUncQ), 2026-11-13; not yet checked by eye.
\scanpage{119/78}
\margin{\unsure{45,49,53}}
\margin{\unsure{33,3\illegible{1}},26,8,}
\margin{1,2,9,10}
\margin{\unsure{(17)}}

\uline{ESQUISSE D'UN PROGRAMME}

par Alexandre Grothendieck

\uline{Sommaire} :

1. Envoi.

2. Un ``jeu de Légo-Teichmüller'', et le groupe de Galois de $\overline{\mathbb{Q}}$ sur $\mathbb{Q}$.

3. Corps de nombres associés à un dessin d'enfant.

4. Polyèdres réguliers et corps finis.

5. Haro sur la topologie dite ``générale'', et réflexions heuristiques
vers une topologie dite ``modérée''.

6. ``Théories différentiables'' (à la Nash) et ``théories modérées''.

7. A la Poursuite des Champs.

8. Digressions de géométrie bidimensionelle.

9. Bilan d'une activité enseignante.

10. Epilogue.

Notes.

N.B. Les astéristiques (*) renvoient aux notes figurant au bas de
la même page, les renvois numérotés de (\textsuperscript{1}) à (\textsuperscript{7}) aux notes (rajou-%
tées ultérieurement) réunies à la fin du rapport.

\mspage{\hand{78}}

\clearpage
% Esquisse d'un programme (cote 119, copy 2): archive cote 119, site page 79 (PDF page 80), typescript page 1.
% typescript. Draft compared: esquisse/pages/esq_119_003.tex (copy 1, site 3).
% AI-transcribed by Claude Opus 5.5 (claude-opus-5-5, adaptive thinking, effort high), Message Batches API, batch pilotnext (msgbatch_01A9sdu8SYq23bEtqX9tUncQ), 2026-11-13; not yet checked by eye.
\scanpage{119/79}
\mspage{1}

\deleted{ESQUISSE D'UN PROGRAMME}

\deleted{par A. Grothendieck}

1. Comme la conjoncture actuelle rend de plus en plus illu-%
soires pour moi les perspectives d'un enseignement de recherche à
l'Université, je me suis résolu à demander mon admission au CNRS, pour
pouvoir consacrer mon énergie à développer des travaux et perspectives
dont il devient clair qu'il ne se trouvera aucun élève (ni même,
semble-t-il, aucun congénère mathématicien) pour les développer à ma
place.

En guise de document ``Titres et Travaux'', on trouvera à la
suite de ce texte la reproduction intégrale d'une esquisse, par thèmes,
de ce que je considérais comme mes principales contributions mathéma-%
tiques au moment d'écrire ce rapport, en 1972. Il contient également
une liste d'articles publiés à cette date. J'ai cessé toute publica-%
tion d'articles scientifiques depuis 1970. Dans les lignes qui suivent,
je me propose de donner un aperçu au moins sur quelques thèmes princi-%
paux de mes réflexions mathématiques depuis lors. Ces réflexions se
sont matérialisées au cours des années en deux volumineux cartons de
notes manuscrites, difficilement déchiffrables sans doute à tout autre
qu'à moi-même, et qui, après des stades de décantations successives,
attendent leur heure peut-être pour une rédaction d'ensemble tout au
moins provisoire, à l'intention de la communauté mathématique. Le
terme ``rédaction'' ici est quelque peu impropre, alors qu'il s'agit
bien plus de développer des idées et visions multiples amorcées au
cours de ces douze dernières années, en les précisant et les appro-%
fondissant, avec tous les rebondissements imprévus qui constamment
accompagnent ce genre de travail - un travail de découverte donc, et
non de compilation de notes pieusement accumulées. Et je compte bien,
dans l'écriture des ``Réflexions Mathématiques'' commencée depuis fé-%
vrier 1983, laisser apparaître clairement au fil des pages la démarche
de la pensée qui sonde et qui découvre, en tâtonnant dans la pénombre
bien souvent, avec des trouées de lumière subites quand quelque tenace

\margin{79}

\clearpage
% Esquisse d'un programme (cote 119, copy 2): archive cote 119, site page 80 (PDF page 81), typescript page 2.
% typescript. Draft compared: esquisse/pages/esq_119_004.tex (copy 1, site 4).
% AI-transcribed by Claude Opus 5.5 (claude-opus-5-5, adaptive thinking, effort high), Message Batches API, batch esq119_full (msgbatch_015xqDHkTumDysRF1qngZhRJ), 2026-11-13; not yet checked by eye.
\scanpage{119/80}
\mspage{2}

image fausse, ou simplement inadéqua\deleted{n}te, se trouve enfin débusquée
et mise à jour, et que les choses qui semblaient de guingois se
mettent en place, dans l'harmonie mutuelle qui leur est propre.
\qquad Quoi qu'il en soit, l'esquisse qui suit de quelques
thèmes de réflexion des dernières dix ou douze années, tiendra
lieu en même temps d'esquisse de programme de travail pour les
années qui viennent, que je compte consacrer au développement de
ces thèmes, ou au moins de certains d'entre eux. Elle est desti-%
née, d'une part aux collègues du Comité National appelés à statuer
sur ma demande, d'autre part à quelques autres collègues, anciens
élèves, amis, dans l'éventualité où certaines des idées esquissées
ici pourraient intéresser l'un d'entre eux.

2. \qquad Les exigences d'un enseignement universitaire, s'adressant
donc à des étudiants (y compris les étudiants dits ``avancés'') au ba-%
gage mathématique modeste (et souvent moins que modeste), m'ont amené
à renouveler de façon draconienne les thèmes de réflexion à proposer
à mes élèves, et de fil en aiguille et de plus en plus, à moi-même
également. Il m'avait semblé important de partir d'un bagage intuitif
commun, indépendant de tout langage technique censé l'exprimer, bien
antérieur à tout tel langage - il s'est avéré que l'intuition géomé-%
trique et topologique des formes, et plus particulièrement des formes
bidimensionelles, était un tel terrain commun. Il s'agit donc de thèmes
qu'on peut grouper sous l'appellation de ``topologie des surfaces'' ou
``géométrie des surfaces'', étant entendu dans cette dernière appella-%
tion que l'accent principal se trouve sur les propriétés topologiques
des surfaces, ou sur les aspects combinatoires qui en constituent
l'expression technique la plus terre-à-terre, et non sur les aspects
différentiels, voire conformes, riemaniens, holomorphes et (de là)
l'aspect ``courbes algébriques complexes''. Une fois ce dernier pas
franchi cependant, voici soudain la géométrie algébrique (mes anciennes
amours !) qui fait irruption à nouveau, et ce par les objets qu'on
peut considérer comme les pierres de construction ultimes de toutes
les autres variétés algébriques. Alors que dans mes recherches d'avant

\margin{80}

\clearpage
% Esquisse d'un programme (cote 119, copy 2): archive cote 119, site page 81 (PDF page 82), typescript page 3.
% typescript. Draft compared: esquisse/pages/esq_119_005.tex (copy 1, site 5).
% AI-transcribed by Claude Opus 5.5 (claude-opus-5-5, adaptive thinking, effort high), Message Batches API, batch esq119_full (msgbatch_015xqDHkTumDysRF1qngZhRJ), 2026-11-13; not yet checked by eye.
\scanpage{119/81}
\mspage{3}

1970, mon attention systématiquement était dirigée vers les objets
de généralité maximale, afin de dégager un langage d'ensemble adéquat
pour le monde de la géométrie algébrique, et que je ne m'attardais
sur les courbes algébriques que dans la stricte mesure où cela s'avé-%
rait indispensable (notamment en cohomologie étale) pour développer
des techniques et énoncés ``passe-partout'' valables en toutes dimen-%
sions et en tous lieux (j'entends, sur tous schémas de base, voire
tous topos annelés de base ...), me voici donc ramené, par le tru-%
chement d'objets si simples qu'un enfant peut les connaître en jouant,
aux débuts et origines de la géométrie algébrique, familiers à Riemann
et à ses émules !

Depuis environ 1975, c'est donc la géométrie des surfaces
(réelles), et à partir de 1977 les liens entre les questions de géo-%
métrie des surfaces et la géométrie algébrique des courbes algébri-%
ques définies sur des corps tels que $\mathbb{C}$, $\mathbb{R}$ ou des extensions de type
fini de $\mathbb{Q}$, qui ont été ma principale source d'inspiration, ainsi
que mon fil conducteur constant. C'est avec surprise et avec émer-%
veillement qu'au fil des ans je découvrais (ou plutôt, sans doute,
redécouvrais) la richesse prodigieuse, réellement inépuisable, la
profondeur insoupçonnée de ce thème, d'apparence si anodine. Je crois
y sentir un point névralgique entre tous, un point de convergence
privilégié des principaux courants d'idées mathématiques, comme aussi
des principales structures et des visions des choses qu'elles expri-%
ment, depuis les plus spécifiques (tels les anneaux $\mathbb{Z}$, $\mathbb{Q}$, $\bar{\mathbb{Q}}$, $\mathbb{R}$, $\mathbb{C}$,
ou le groupe $S\hand{\ell}(2)$ sur l'un de ces anneaux, ou les groupes algébri-%
ques réductifs généraux) aux plus ``abstraits'', telles les ``multipli-%
cités'' algébriques, analytiques complexes ou analytiques réelles.
(Celles-ci s'introduisent naturellement quand il s'agit d'étudier
systématiquement des ``variétés de modules'' pour les objets géométri-%
ques envisagés, si on veut dépasser le point de vue notoirement in-%
suffisant des ``modules grossiers'', qui revient à ``tuer'' bien malen-%
contreusement les groupes d'automorphismes de ces objets .) Parmi ces
multiplicités modulaires, ce sont celles de Mumford-Deligne pour les
courbes algébriques ``stables'' de genre $g$, à $\nu$ points marqués, que

\margin{81}

\clearpage
% Esquisse d'un programme (cote 119, copy 2): archive cote 119, site page 82 (PDF page 83), typescript page 4.
% typescript. Draft compared: esquisse/pages/esq_119_006.tex (copy 1, site 6).
% AI-transcribed by Claude Opus 5.5 (claude-opus-5-5, adaptive thinking, effort high), Message Batches API, batch esq119_full (msgbatch_015xqDHkTumDysRF1qngZhRJ), 2026-11-13; not yet checked by eye.
\scanpage{119/82}
\mspage{4}

je note $M^{\wedge}_{g,\nu}$ (compactification de la multiplicité ``ouverte''
$M_{g,\nu}$ correspondant aux courbes lisses), qui depuis quelques deux
ou trois années ont exercé sur moi une fascination particulière,
plus forte peut-être qu'aucun autre objet mathématique à ce jour.
A vrai dire, il s'agit plutôt du système de \uline{toutes} les multiplicités
$M_{g,\nu}$ pour $g,\nu$ variables, liées entre elles par un certain nombre
d'opérations fondamentales (telles les opérations de ``bouchage de
trous'' i.e. de ``gommage'' de points marqués, celle de ``recollement'',
et les opérations inverses), qui sont le reflet en géométrie algé-%
brique absolue de caractéristique zéro (pour le moment) d'opérations
géométriques familières du point de vue de la ``chirurgie'' topologique
ou conforme des surface. La principale raison sans doute de cette
fascination, c'est que cette structure géométrique très riche sur le
système des multiplicités modulaires ``ouvertes'' $M_{g,\nu}$ se réflète
par une structure analogue sur les groupoïdes fondamentaux corres-%
pondants, les ``groupoïdes de Teichmüller'' $\hat{T}_{g,\nu}$ ,et que ces opérations
au niveau des $\hat{T}_{g,\nu}$ ont un caractère suffisamment intrinsèque pour
que le groupe de Galois $\mathrm{I}\!\Gamma$ de $\overline{\mathbb{Q}}/\mathbb{Q}$ opère sur toute cette ``tour'' de
groupoïdes de Teichmüller, en respectant toutes ces structures. Chose
plus extraordinaire encore, cette opération est \uline{fidèle} - à vrai dire,
elle est fidèle déjà sur le premier ``étage'' non trivial de cette tour,
savoir $\hat{T}_{0,4}$ - ce qui signifie aussi, essentiellement, que l'action
extérieure de $\mathrm{I}\!\Gamma$ sur le groupe fondamental $\hat{\pi}_{0,3}$ de la droite pro-%
jective standard $\mathbb{P}^1$ sur $\mathbb{Q}$, privée des trois points $0,\ 1,\ \infty$, est
déjà fidèle. Ainsi \uline{le groupe de Galois} $\mathrm{I}\!\Gamma$ \uline{se réalise comme un groupe}
\uline{d'automorphisme d'un groupe profini des plus concrets}, respectant
d'ailleurs certaines structures essentielles de ce groupe. Il s'ensuit
qu'un élément de $\mathrm{I}\!\Gamma$ peut être ``paramétré'' (de diverses façons équi-%
valentes d'ailleurs) par un élément convenable de ce groupe profini
$\hat{\pi}_{0,3}$ (un groupe profini libre à deux générateurs), ou par un système
de tels éléments, ce ou ces éléments étant d'ailleurs soumis à cer-%
taines conditions simples, nécessaires (et sans doute non suffisantes)
pour que ce ou ces éléments corresponde(nt) bien à un élément de $\mathrm{I}\!\Gamma$.
Une des tâches les plus fascinantes ici, est justement d'appréhender
des conditions nécessaires \uline{et} suffisantes sur un automorphisme extérieur

\margin{82}

\clearpage
% Esquisse d'un programme (cote 119, copy 2): archive cote 119, site page 83 (PDF page 84), typescript page 5.
% typescript. Draft compared: esquisse/pages/esq_119_007.tex (copy 1, site 7).
% AI-transcribed by Claude Opus 5.5 (claude-opus-5-5, adaptive thinking, effort high), Message Batches API, batch esq119_full (msgbatch_015xqDHkTumDysRF1qngZhRJ), 2026-11-13; not yet checked by eye.
\scanpage{119/83}
\mspage{5}

de $\hat{\pi}_{0,3}$ i.e. sur le ou les paramètres correspondants, pour qu'il
provienne d'un élément de $I\!\Gamma$ - ce qui fournirait une description
``purement algébrique'', en termes de groupes profinis et sans ré-%
férence à la théorie de Galois des corps de nombres, du groupe de
Galois $I\!\Gamma = \mathrm{Gal}(\bar{\underline{Q}}/\underline{Q})$ !

Peut-être une caractérisation même conjecturale de $I\!\Gamma$
comme sous-groupe de Autext($\hat{\pi}_{0,3}$) est-elle pour le moment hors de
portée($^1$) ; je n'ai pas de conjecture à proposer encore. Une autre
tâche par contre est abordable immédiatement, c'est celle de décrire
l'action de $I\!\Gamma$ sur toute la tour de Teichmüller, en termes de son
action sur le ``premier étage'' $\hat{\pi}_{0,3}$ , i.e. exprimer un automorphisme
de cette tour, en termes du ``paramètre'' dans $\hat{\pi}_{0,3}$ qui repère l'élé-%
ment courant $\gamma$ de $I\!\Gamma$. Ceci est lié à une représentation de la tour
de Teichmüller (en tant que groupoïde muni d'une opération de ``recol-%
lement'') par générateurs et relations, qui donnera en particulier
une présentation par générateurs et relations, au sens ordinaire, de
chacun des $\hat{T}_{g,\nu}$ (en tant que groupoïde profini). Ici, même pour
$g = 0$ (donc quand les groupes de Teichmüller correspondants sont
des groupes de tresses ``bien connus''), les générateurs et relations
connus à ce jour dont j'ai eu connaissance, me semblent inutilisables
tels quels, car ils ne présentent pas les caractères d'invariance et
de symétrie indispensables pour que l'action de $I\!\Gamma$ soit directement
lisible sur cette présentation. Ceci est lié notamment au fait que
les gens s'obstinent encore, en calculant avec des groupes fondamen-%
taux, à fixer un seul point base, plutôt que d'en choisir astucieu-%
sement tout un paquet qui soit invariant par les symétries de la si-%
tuation, lesquelles sont donc perdues en route. Dans certaines situa-%
tions (comme des théorèmes de descente à la Van Kampen pour groupes
fondamentaux) il est bien plus élégant, voir indispensable pour y com-%
prendre quelque chose, de travailler avec des groupoïdes fondamentaux
par rapport à un paquet de points bases convenable , et il en est cer-%
tainement ainsi pour la tour de Teichmüller. Il semblerait (incroyable,
mais vrai !) que la géométrie même du premier étage de la tour de
Teichmüller (correspondant donc aux ``modules'' soit pour des droites

\margin{83}

\clearpage
% Esquisse d'un programme (cote 119, copy 2): archive cote 119, site page 84 (PDF page 85), typescript page 6.
% typescript. Draft compared: esquisse/pages/esq_119_008.tex (copy 1, site 8).
% AI-transcribed by Claude Opus 5.5 (claude-opus-5-5, adaptive thinking, effort high), Message Batches API, batch esq119_full (msgbatch_015xqDHkTumDysRF1qngZhRJ), 2026-11-13; not yet checked by eye.
\scanpage{119/84}
\mspage{6}

projectives avec quatre points marqués, soit pour des courbes ellip-%
tiques (!)) n'ait jamais été bien explicité, par exemple la relation
entre le cas de genre 0 avec la géométrie de l'octaèdre, et celle du
tétraèdre. A fortiori les multiplicités modulaires $M_{0,5}$ (pour les
droites projectives avec cinq points marqués) et $M_{1,2}$ (pour les
courbes de genre 1 avec deux points marqués), d'ailleurs quasiment
isomorphes entre elles, semblent-elles terre vierge - les groupes
de tresses ne vont pas nous éclairer à leur sujet ! J'ai commencé à
regarder $M_{0,5}$ à des moments perdus, c'est un véritable joyau, d'une
géométrie très riche étroitement liée à celle de l'icosaèdre.

L'intérêt a priori d'une connaissance complète des deux
premiers étages de la tour (savoir, les cas où la dimension modulaire
$N = 3g-3+\nu$ est $\leq 2$) réside dans ce principe, que \uline{la tour entière se
reconstitue à partir des deux premiers étages}, en ce sens que via
l'opération fondamentale de ``recollement'', l'étage 1 fournit un sys-%
tème complet de générateurs, et l'étage 2 un système complet de re-%
lations. Il y a une analogie frappante, et j'en suis persuadé, pas
seulement formelle, entre ce principe, et le principe analogue de
Demazure pour la structure des groupes algébriques réductifs, si on
remplace le terme ``étage'' ou ``dimension modulaire'' par ``rang semi-%
simple du groupe réductif''. Le lien devient plus frappant encore,
si on se rappelle que le groupe de Teichmüller $T_{1,1}$ (dans le contexte
discret transcendant maintenant, et non dans le contexte algébrique
profini, où on trouve les complétions profinies des premiers) n'est
autre que $S\ell(2,\hand{\mathbb{Z}})$, i.e. le groupe des points entiers du schéma en
groupes simple de rang 1 ``absolu'' $S\ell(2)_{\hand{\mathbb{Z}}}$. Ainsi, \uline{la pierre de cons-%
truction fondamentale pour la tour de Teichmüller}, \uline{est essentiellement
la même que celle pour ``la tour'' des groupes réductifs de tous rangs} -
un groupe d'ailleurs dont on peut dire sans doute qu'il est présent
dans toutes les disciplines essentielles des mathématiques.

Ce principe de construction de la tour de Teichmüller n'est
pas démontré à l'heure actuelle - mais je n'ai aucun doute qu'il ne
soit valable. Il résulterait (via une théorie de dévissage des struc-%
tures stratifiées - en l'occurence les $M^{\wedge}_{g,\nu}$ - qui resterait à écrire,
cf. par. 5) d'une propriété extrêmement plausible des multiplicités

\margin{84}

\clearpage
% Esquisse d'un programme (cote 119, copy 2): archive cote 119, site page 85 (PDF page 86), typescript page 7.
% typescript. Draft compared: esquisse/pages/esq_119_009.tex (copy 1, site 9).
% AI-transcribed by Claude Opus 5.5 (claude-opus-5-5, adaptive thinking, effort high), Message Batches API, batch esq119_full (msgbatch_015xqDHkTumDysRF1qngZhRJ), 2026-11-13; not yet checked by eye.
\scanpage{119/85}
\mspage{7}

modulaires ouvertes $M_{g,\nu}$ dans le contexte analytique complexe,
à savoir que pour une dimension modulaire $N \geq 3$, le groupe fonda-%
mental de $M_{g,\nu}$ (i.e. le groupe de Teichmüller habituel $T_{g,\nu}$)
est isomorphe au ``groupe fondamental à l'infini'' i.e. celui d'un
``voisinage tubulaire de l'infini''. C'est là une chose bien familière
(due à Lefschetz essentiellement) pour une variété lisse \uline{affine} de
dimension $N \geq 3$. Il est vrai que les multiplicités modulaires ne
sont pas affines (sauf pour des petites valeurs de $g$), mais il suf-%
firait qu'une telle $M_{g,\nu}$ de dimension $N$ (ou plutôt, un revête-%
\margin{\unsure{$\frown$}}ment fini convena\deleted{n}ble) soit réunion de $N$-2 ouverts affines, donc
que $M_{g,\nu}$ ne soit pas ``trop proche d'une variété compacte''.

N'ayant aucun doute sur ce principe de construction de la
tour de Teichmüller, je préfère laisser aux experts de la théorie
transcendante, mieux outillés que moi, le soin de prouver le néces-%
saire (s'il s'en trouve qui soit intéressé), pour expliciter plutôt,
avec tout le soin qu'elle mérite, la structure qui en découle pour
la tour de Teichmüller par générateurs et relations, dans le cadre
discret cette fois et non profini - ce qui revient, essentiellement,
à une compréhension complète des quatre multiplicités modulaires
$M_{0,4}$, $M_{1,1}$, $M_{0,5}$, $M_{1,2}$ , et de leurs groupoïdes fondamentaux par rap-%
port à des ``points base'' convenablement choisis. Ceux-ci s'offrent
tout naturellement, comme les courbes algébriques complexes du type
$(g,\nu)$ envisagé, qui ont un groupe d'automorphismes (nécessairement
fini) plus grand que dans le cas générique (*). En y incluant la
sphère holomorphe à trois points marqués (provenant de $M_{0,3}$ i.e.

\rule{4cm}{0.4pt}

(*) Il faut y ajouter de plus les ``points-base'' provenant par opéra-%
tions de recollement de ``pièces'' du même type en dimension modulaire
inférieure. D'autre part, en dimension modulaire 2 (cas de $M_{0,5}$ et $M_{1,2}$),
il convient d'exclure les points de certaines familles à un paramètre
de courbes admettant un automorphisme exceptionnel d'ordre 2. Ces fa-%
milles constituent d'ailleurs sur les multiplicités envisagées des cour-%
bes rationnelles remarquables, qui me paraissent un ingrédient important
de la structure de ces multiplicités.

\margin{85}

\clearpage
% Esquisse d'un programme (cote 119, copy 2): archive cote 119, site page 86 (PDF page 87), typescript page 8.
% typescript. Draft compared: esquisse/pages/esq_119_010.tex (copy 1, site 10).
% AI-transcribed by Claude Opus 5.5 (claude-opus-5-5, adaptive thinking, effort high), Message Batches API, batch esq119_full (msgbatch_015xqDHkTumDysRF1qngZhRJ), 2026-11-13; not yet checked by eye.
\scanpage{119/86}
\mspage{8}

de l'étage O), on trouve \uline{douze} ``\uline{pièces de construction}'' \uline{fondamentales}
(6 de genre O, 6 de genre 1) dans un ``jeu de Légo-Teichmüller\hand{\unsure{*}}(grande \margin{\unsure{*}}
boite), où les points marqués sur les surfaces envisagées sont rem-%
placés par des ``trous'' à bord, de façon à avoir des surfaces à bord,
donc des pièces de construction qui peuvent s'assembler par frotte-%
ment doux comme dans le jeu de Légo ordinaire cher à nos enfants
(ou petits-enfants...). Par assemblage on trouve un moyen tout ce
qu'il y a de visuel pour construire tout type de surface (ce sont
ces assemblages essentiellement qui seront les ``points bases'' pour
notre fameuse tour), et aussi de visualiser les ``\uline{chemins}'' \uline{élémentaires}
par des opérations tout aussi concrètes tels des ``twists'', ou des
automorphismes des pièces du jeu, et d'écrire les \uline{relations fondamentales}
entre chemins composés. Suivant la taille (et le prix !) de la boite
de construction utilisée, on \unsure{trcuve|trouve} d'ailleurs de nombreuses descrip-%
tions différentes de la tour de Teichmüller par générateurs et rela-%
tions. La boite la plus petite est réduite à des pièces toutes iden-%
tiques, de type (O,3) - ce sont les ``pantalons'' de Thurston, et le
jeu de Légo-Teichmüller que j'essaie de décrire, issu de motivations
et de réflexions de géométrie algébrique absolue sur le corps \hand{$\underline{Q}$},
est très proche du jeu de ``chirurgie géodésique hyperbolique'' de
Thurston, dont j'ai appris l'existence l'an dernier par Yves Ladegaillerie.
Dans un microséminaire avec Carlos Contou-Carrère et Yves Ladegraillerie,
nous avons amorcé une réflexion dont un des objets est de confronter
les deux points de vue, qui se complètent mutuellement.

J'ajoute que chacune des douze pièces de construction de
la ``grande boite'' se trouve muni d'une décomposition cellulaire cano-%
nique, stable par toutes les symétries, ayant comme seuls sommets les
``points marqués'' (ou centres des trous), et comme arètes certains che-%
mins géodésiques (pour la structure riemanienne canonique sur la sphère
ou le tore envisagé) entre certaines paires de sommets (savoir ceux
qui se trouvent sur un même ``lieu réel'', pour une structure réelle
convenable de la courbe algébrique complexe envisagée). Par suite,\inserted{dans ce jeu}
% NOTE: by hand, "dans ce jeu" in the next line is circled and joined by a line to the place after "Par suite," (a move mark); the words are not struck through
toutes les surfaces obtenues dans ce jeu par assemblage sont munies
de structures cellulaires canoniques, qui à leur tour (cf.§3 plus bas)

\margin{86}

\clearpage
% Esquisse d'un programme (cote 119, copy 2): archive cote 119, site page 87 (PDF page 88), typescript page 9.
% typescript. Draft compared: esquisse/pages/esq_119_011.tex (copy 1, site 11).
% AI-transcribed by Claude Opus 5.5 (claude-opus-5-5, adaptive thinking, effort high), Message Batches API, batch esq119_full (msgbatch_015xqDHkTumDysRF1qngZhRJ), 2026-11-13; not yet checked by eye.
\scanpage{119/87}
\mspage{9}

permettent de considérer ces surfaces comme associées à des courbes
algébriques complexes (et même sur $\overline{\mathbb{Q}}$) canoniquement déterminées.
Il y a là un jeu de chassé-croisé typique entre le combinatoire,
et l'algébrique complexe (ou mieux, l'algébrique sur $\overline{\mathbb{Q}}$).

La ``petite boite'' aux pièces toutes identiques, qui a le
charme de l'économie, donnera sans doute une description relative-%
ment compliquée pour les relations (compliquée, mais nullement inex-%
tricable !). La grande boite donnera lieu à des relations plus nom-%
breuses (du fait qu'il y a beaucoup plus de points-bases et de che-%
mins remarquables entre eux), mais à structure plus transparente. Je
prévois qu'en dimension modulaire 2, tout comme dans le cas plus ou
moins familier de la dimension modulaire 1 (avec notamment la des-%
\margin{//} cription de $S\ell(2,\mathbb{Z})$ par $(\rho,\sigma \mid \rho^3 = \sigma^2,\ \over{\rho}{\sigma}^4 = \over{\sigma}{\rho}^6 = 1)$), on trouvera
un engendrement par les groupes d'automorphismes des trois types de
pièces pertinentes, avec des relations simples que je n'ai pas déga-%
gé à l'heure d'écrire ces lignes. Peut-être même trouvera-t-on un
principe de ce genre pour tous les $T_{g,\nu}$, ainsi qu'une décomposition
cellulaire de $M^{\wedge}_{g,\nu}$ généralisant celles qui se présentent sponta-%
nément pour $M^{\wedge}_{0,4}$ et $M^{\wedge}_{1,1}$, et que j'entrevois dès à présent pour
la dimension modulaire 2, en utilisant les hypersurfaces correspon-%
dant aux diverses \uline{structures réelles} sur les structures complexes
envisagées, pour effectuer le découpage cellulaire voulu.


3. Plutôt que de suivre (comme prévu) un ordre thématique ri-%
goureux, je me suis laissé emporter par ma prédilection pour un thème
particulièrement riche et brûlant, auquel je compte me consacrer
d'ailleurs prioritairement pendant quelque temps, à partir de la ren-%
trée 84/85. Je reprends donc l'exposé thématique là où je l'ai laissé,
tout au début du paragraphe précédent.

Mon intérêt pour les surfaces topologiques commence à poindre
en 1974, où je propose à Yv\unsure{è|e}s Ladegaillerie le thème de l'étude iso-%
topique des plongements d'un 1-complexe topologique dans une surface
compacte. Dans les deux années qui suivent, cette étude le conduit à
un remarquable théorème d'isotopie, donnant une description algébrique

\margin{87}

\clearpage
% Esquisse d'un programme (cote 119, copy 2): archive cote 119, site page 88 (PDF page 89), typescript page 10.
% typescript. Draft compared: esquisse/pages/esq_119_012.tex (copy 1, site 12).
% AI-transcribed by Claude Opus 5.5 (claude-opus-5-5, adaptive thinking, effort high), Message Batches API, batch esq119_full (msgbatch_015xqDHkTumDysRF1qngZhRJ), 2026-11-13; not yet checked by eye.
\scanpage{119/88}
\mspage{10}

complète des classes d'isotopie de plongements de tels 1-complexes,
ou de surfaces compactes à bord, dans une surface compacte orientée,
en termes de certains invariants combinatoires très simples, et des
groupes fondamentaux des protagonistes. Ce théorème, qui doit pou-%
voir s'étendre sans mal aux plongements d'un espace compact quelcon-%
que (triangulable pour simplifier) dans une surface compacte orientée,
redonne comme corollaires faciles plusieurs résultats classiques
profonds de théorie des surfaces, et notamment le théorème d'isoto-%
pie de Baer. Il va finalement être publié, séparément du reste (et
dix ans après, vu la dureté des temps...), dans Topology. Dans le
travail de Ladegaillerie figure également une description purement
algébrique, en termes de groupoïdes fondamentaux, de la catégorie
``isotopique'' des surfaces compactes $X$, munies d'un 1-complexe to-%
pologique $K$ plongé dans $X$. Cette description, qui a eu le malheur
d'aller à l'encontre du ``goût du jour'' et de ce fait semble i\over{m}{n}publia-% \margin{n}
ble, a néanmoins servi (et sert encore) comme un guide précieux dans
mes réflexions ultérieures, notamment dans le contexte de la géomé-%
trie algébrique absolue de caractéristique nulle.

Le cas où $(X,K)$ est une ``carte'' 2-dimensionnelle, i.e.
où les composantes connexes de $X\backslash K$ sont des 2-cellules ouvertes
(et où de plus $K$ est muni d'un ensemble fini $S$ de ``sommets'',
tel que les composantes connexes de $K\backslash S$ soient des 1-cellules ou-%
vertes) attire progressivement mon attention dans les années suivantes.
La catégorie isotopique de ces cartes admet une description algébrique
particulièrement simple, via l'ensemble des ``repères'' (ou ``drapeaux'',
ou ``biarcs'') associés à la carte, qui se trouve naturellement muni
d'une structure d'ensemble à groupe d'opérateurs, sous le groupe
\[ \underline{C}_2 = \langle \sigma_0,\sigma_1,\sigma_2 \,|\, \sigma_0^{\,2} = \sigma_1^{\,2} = \sigma_2^{\,2} = (\sigma_0\sigma_2)^2 = 1 \rangle \ , \]
que j'appelle le \uline{groupe cartographique} (non orienté) de dimension 2.
Iladmet comme sous-groupe d'indice 2 le \uline{groupe cartographique orienté},
engendré par les produits en nombre pair des générateurs, qui peut
aussi se décrire comme
\[ \underline{C}_2^{+} = \langle \rho_s,\rho_f,\sigma \,|\, \rho_s\rho_f = \sigma,\ \sigma^2 = 1 \rangle \ , \]

\margin{88}

\clearpage
% Esquisse d'un programme (cote 119, copy 2): archive cote 119, site page 89 (PDF page 90), typescript page 11.
% typescript. Draft compared: esquisse/pages/esq_119_013.tex (copy 1, site 13).
% AI-transcribed by Claude Opus 5.5 (claude-opus-5-5, adaptive thinking, effort high), Message Batches API, batch esq119_full (msgbatch_015xqDHkTumDysRF1qngZhRJ), 2026-11-13; not yet checked by eye.
\scanpage{119/89}
\mspage{11}

(avec
\[ \rho_s = \sigma_2\sigma_1 \ , \quad \rho_f = \sigma_1\sigma_0 \ , \quad \sigma = \sigma_0\sigma_2 = \sigma_2\sigma_0 \ , \]
opérations de \uline{rotation élémentaire} d'un repère autour d'un sommet,
d'une face et d'une arête respectivement). Il y a un dictionnaire
parfait entre la situation topologique des cartes compactes, resp.
cartes compactes orientées, d'une part, et les ensembles finis à
groupe d'opérateurs $\underline{C}_2$ resp. $\underline{C}_2^+$ de l'autre, dictionnaire dont
l'existence était d'ailleurs plus ou moins connue, mais jamais
énoncée avec la précision nécessaire, ni développé tant soit peu.
Ce travail de fondements est fait avec le soin qu'il mérite dans
un excellent travail de DEA, fait en commun par Jean Malgoire et
Christine Voisin en 1976.

Cette réflexion prend soudain une dimension nouvelle,
avec cette remarque simple que le groupe $\underline{C}_2^+$ peut s'interpréter
comme un quotient du groupe fondamental d'une sphère orientée pri-%
vée de trois points, numérotés 0, 1, 2, les opérations $\rho_s$, $\sigma$, $\rho_f$
s'interprétant comme les lacets autour de ces points, satisfaisant
la relation familière
\[ \hand{\ell}_0\hand{\ell}_1\hand{\ell}_2 = 1 \ , \]
alors que la relation supplémentaire $\sigma^2 = 1$ i.e. $\hand{\ell}_1^{\,2} = 1$ signi-%
fie qu'on s'intéresse au quotient du groupe fondamental correspon-%
dant à un indice de ramification imposé 2 au point 1, qui classi-%
fie donc les revêtements de la sphère, ramifiés au plus en les
points 0, 1, 2, avec une ramification égale à 1 ou 2 en les points
au-dessus de 1. Ainsi, les cartes orientées compactes forment une
catégorie isotopique équivalente à celle de ces revêtements, soumis
de plus à la condition supplémentaire d'être des revêtements finis.
Prenant maintenant comme sphère de référence la sphère de Riemann,
ou droite projective complexe, rigidifiée par les trois point 0, 1
et $\infty$ (ce dernier remplaçant donc 2), et se rappelant que tout revê-%
tement ramifié fini d'une courbe algébrique complexe hérite lui-même
d'une structure de courbe algébrique complexe, on aboutit à cette
constatation, qui huit ans après me paraît encore toujours aussi
extraordinaire : \uline{toute carte orientée} ``\uline{finie}'' \uline{se réalise canoniquement}
\uline{sur une courbe algébrique complexe} ! Mieux encore, comme la droite

\margin{89}

\clearpage
% Esquisse d'un programme (cote 119, copy 2): archive cote 119, site page 90 (PDF page 91), typescript page 12.
% typescript. Draft compared: esquisse/pages/esq_119_014.tex (copy 1, site 14).
% AI-transcribed by Claude Opus 5.5 (claude-opus-5-5, adaptive thinking, effort high), Message Batches API, batch esq119_full (msgbatch_015xqDHkTumDysRF1qngZhRJ), 2026-11-13; not yet checked by eye.
\scanpage{119/90}
\mspage{12}

projective complexe est définie sur le corps de base absolue $\mathbb{Q}$ ,
ainsi que les points de ramification admis, les courbes algébriques
obtenues sont définies non seulement sur $\mathbb{C}$, mais sur la clôture
algébrique $\overline{\mathbb{Q}}$ de $\mathbb{Q}$ dans $\mathbb{C}$. Quant à la carte de départ, elle se
retrouve sur la courbe algébrique, comme image inverse du segment
réel $[0,1]$ (où $0$ est considéré comme un sommet, et $1$ comme
milieu d'une ``arête pliée'' ayant 1 comme centre), lequel constitue
dans la sphère de Riemann la ``2-carte orientée universelle'' (*).
Les points de la courbe algébrique $X$ au-dessus de 0, de 1 et de \unsure{\hand{$\infty$}|$\infty$}
ne sont autres que les sommets, et les ``centres'' des arêtes et des
faces respectivement de la carte $(X,K)$, et les ordres des sommets
et des faces ne sont autres que les multiplicités des zéros et
des pôles de la fonction rationnelle (définie sur $\overline{\mathbb{Q}}$) sur $X$,
exprimant sa projection structurale vers $\mathbb{P}^1_{\mathbb{C}}$.

Cette découverte, qui techniquement se réduit à si peu
de choses, a fait sur moi une impression très forte, et elle repré-%
sente un tournant décisif dans le cours de mes réflexions, un dé-%
placement notamment de mon centre d'intérêt en mathématique, qui
soudain s'est \inserted{re}trouvé fortement localisé. Je ne crois pas qu'un
fait mathématique m'ait jamais autant frappé que celui-là, et ait
eu un impact psychologique comparable ($^2$). Cela tient sûrement à la
nature tellement familière, non technique, des objets considérés,
dont tout dessin d'enfant griffonné sur un bout de papier (pour peu
que le graphisme soit d'un seul tenant) donne un exemple parfaitement
explicite. A un tel dessin se trouvent associés des invariants arith-%
métiques subtils, qui seront chamboulés complètement dès qu'on y
rajoute un trait de plus. S'agissant ici de cartes sphériques, don-%

\noindent\rule{4cm}{0.4pt}

(*) Il y a une description analogue des cartes finies non orientées,
éventuellement avec bord, en termes de courbes algébriques \uline{réelles},
plus précisément de revêtement de $\mathbb{P}^1_{\mathbb{R}}$ ramifié seulement on 0, 1, $\infty$,
la surface à bord associée à un tel revêtement étant $X(\mathbb{C})/\tau$, où $\tau$
est la conjugaison complexe. La carte non orientée ``universelle'' est
ici le disque, ou hémisphère supérieur de la sphère de Riemann, muni
comme précédemment du 1-complexe plongé $K = [0,1]$.

\margin{90}

\clearpage
% Esquisse d'un programme (cote 119, copy 2): archive cote 119, site page 91 (PDF page 92), typescript page 13.
% typescript. Draft compared: esquisse/pages/esq_119_015.tex (copy 1, site 15).
% AI-transcribed by Claude Opus 5.5 (claude-opus-5-5, adaptive thinking, effort high), Message Batches API, batch esq119_full (msgbatch_015xqDHkTumDysRF1qngZhRJ), 2026-11-13; not yet checked by eye.
\scanpage{119/91}
\mspage{13}

nant nécessairement naissance à des courbes de genre zéro (qui ne
fournissent donc pas des ``modules''), on peut dire que la courbe
en question est ``épinglée'' dès qu'on fixe trois de ses points, par
exemple trois sommets de la carte, ou plus généralement trois cen-%
tres de facettes (sommets, arêtes ou faces) - dès lors l'applica-%
tion structurale $f : X \to \hand{$\mathbb{P}$}^{1}_{\hand{$\mathbb{C}$}}$ peut s'interpréter comme une fonction
rationnelle
\[ f(z) = P(z)/Q(z) \hand{$\in$} \hand{$\mathbb{C}$}(z) \]
bien déterminée, quotient de deux polynômes bien déterminés pre-%
miers entre eux avec $Q$ unitaire, satisfaisant à des conditions
algébriques qui traduisent notamment le fait que $f$ soit non
ramifié en dehors des valeurs $0$, $1$, $\infty$, et qui impliquent que les
coefficients de ces polynômes sont des \uline{nombres algébriques} ; donc
leurs zéros sont des nombres algébriques, qui représentent res-%
pectivement les sommets et les centres des faces de la carte en-%
visagée.

Revenant au cas général, les cartes finies s'interpré-%
tant comme des revêtements sur $\bar{\hand{$\mathbb{Q}$}}$ d'une courbe algébrique définie
sur le corps premier \hand{$\mathbb{Q}$} lui-même, il en résulte que le groupe de
Galois \hand{$\mathbb{I}\Gamma$} de $\bar{\hand{$\mathbb{Q}$}}$ sur \hand{$\mathbb{Q}$} opère sur la catégorie de ces cartes de
façon naturelle. Par exemple, l'opération d'un automorphisme
$\gamma \hand{$\in$} \hand{$\mathbb{I}\Gamma$}$ sur une carte sphérique donnée par la fonction rationnelle
ci-dessus, est obtenue en appliquant $\gamma$ aux coefficients des po-%
lynômes $P$, $Q$. Voici donc ce mystérieux groupe \hand{$\mathbb{I}\Gamma$} intervenir comme
agent transformateur sur des formes topologico-combinatoires de la
nature la plus élémentaire qui soit, amenant à se poser des ques-%
tions comme : telles cartes orientées données sont-elles ``conjuguées'',
ou : quels exactement sont les conjuguées de telle carte orientée
donnée ? (Il y en a, visiblement, un nombre fini seulement.)

J'ai traité quelques cas concrets (pour des revêtements
de bas degrés) par des expédients divers, J. Malgoire en a traité
quelques autres - je doute qu'il y ait une méthode uniforme permet-%
tant d'y répondre à coups d'ordinateurs. Ma réflexion très vite

\margin{91}

\clearpage
% Esquisse d'un programme (cote 119, copy 2): archive cote 119, site page 92 (PDF page 93), typescript page 14.
% typescript. Draft compared: esquisse/pages/esq_119_016.tex (copy 1, site 16).
% AI-transcribed by Claude Opus 5.5 (claude-opus-5-5, adaptive thinking, effort high), Message Batches API, batch esq119_full (msgbatch_015xqDHkTumDysRF1qngZhRJ), 2026-11-13; not yet checked by eye.
\scanpage{119/92}
\mspage{14}

s'est engagée dans une directions plus conceptuelle, pour arriver
à appréhender la nature de cette action de $\mathrm{I\!\Gamma}$. On s'aperçoit
d'emblée que grosso-modo cette action est exprimée par une cer-%
taine action ``extérieure'' de $\mathrm{I\!\Gamma}$ sur le compactifié profini du
groupe cartographique orienté $\underline{C}_2^{+}$ , et cette action à son tour
est déduite par passage au quotient de l'action extérieure cano-%
nique de $\mathrm{I\!\Gamma}$ sur le groupe fondamental profini $\hat{\pi}_{0,3}$ de $(U_{0,3})_{\bar{\mathbb{Q}}}$ ,
où $U_{0,3}$ désigne la courbe-type de genre $0$ sur le corps premier
$\mathbb{Q}$, privée de trois points. C'est ainsi que mon attention s'est
portée vers ce que j'ai appelé depuis la ``\uline{géométrie algébrique}
\uline{anabélienne}'', dont le point de départ est justement une étude
(pour le moment limitée à la caractéristique zéro) de l'action
de groupes de Galois ``absolus'' (notamment les groupes $\mathrm{Gal}(\bar{K}/K)$,
où $K$ est une extension de type fini du corps premier) sur des
groupes fondamentaux géométriques (profinis) de variétés algé-%
briques (définies sur $K$), et plus particulièrement (rompant avec
une tradition bien enracinée) des groupes fondamentaux qui sont
très éloignés des groupes abéliens (et que pour cette raison je
nomme ``\uline{anabéliens}''). Parmi ces groupes, et très proche du groupe
$\hat{\pi}_{0,3}$ , il y a le compactifié profini du groupe modulaire $S\ell(2,\mathbb{Z})$,
dont le quotient par le centre $\pm 1$ contient le précédent comme
sous-groupe de congruence mod 2, et peut s'interpréter d'ailleurs
également comme un groupe ``cartographique'' orienté, savoir celui
qui classifie les cartes orientées \uline{triangulées} (i.e. celles dont
les faces sont des triangles ou des monogones).

Toute carte finie orientée donne lieu à une courbe al-%
gébrique projective et lisse définie sur $\bar{\mathbb{Q}}$ , et il se pose alors
immédiatement la question : quelles sont les courbes algébriques
sur $\bar{\mathbb{Q}}$ obtenues ainsi - les obtiendrait-on toutes, qui sait ?
En termes plus savants, serait-il vrai que toute courbe algébrique
projective et lisse définie sur un corps de nombres interviendrait
comme une ``courbe modulaire'' possible pour paramétriser les cour-%
bes elliptiques munies d'une rigidification convenable ? Une telle
supposition avait l'air à tel point dingue que j'étais presque

\margin{92}

\clearpage
% Esquisse d'un programme (cote 119, copy 2): archive cote 119, site page 93 (PDF page 94), typescript page 15.
% typescript. Draft compared: esquisse/pages/esq_119_017.tex (copy 1, site 17).
% AI-transcribed by Claude Opus 5.5 (claude-opus-5-5, adaptive thinking, effort high), Message Batches API, batch esq119_full (msgbatch_015xqDHkTumDysRF1qngZhRJ), 2026-11-13; not yet checked by eye.
\scanpage{119/93}
\mspage{15}

gêné de la soumettre aux compétences en la matière. Deligne consulté
trouvait la supposition dingue en effet, mais sans avoir un contre-%
exemple dans ses manches. Moins d'un an après, au Congrès Interna-%
tional de Helsinki, le mathématicien soviétique Bielyi annonce
justement ce résultat, avec une démonstration d'une simplicité dé-%
concertante tenant en deux petites pages d'une lettre de Deligne -
jamais sans doute un résultat profond et déroutant ne fut démontré
en si peu de lignes !

Sous la forme où l'énonce Bielyi, son résultat dit essen-%
tiellement que \uline{toute courbe algébrique définie sur un corps de nom-}%
\uline{bres peut s'obtenir comme revêtement de la droite projective ramifié}
\uline{seulement en les points} 0, 1, $\infty$. Ce résultat semble être passé plus
ou moins inaperçu. Pourtant, il m'apparaît d'une portée considéra-%
ble. Pour moi, son message essentiel a été qu'\uline{il y a une identité}
\uline{profonde entre la combinatoire des cartes finies d'une part}, \uline{et la}
\uline{géométrie des courbes algébriques définies sur des corps de nombres},
\uline{de l'autre}. Ce résultat profond, joint à l'interprétation algébrico-%
géométrique des cartes finies, ouvre la porte sur un monde nouveau,
inexploré - et à portée de main de tous, qui passent sans le voir.

C'est près de trois ans plus tard seulement, voyant que
décidément les vastes horizons qui s'ouvrent là ne faisaient rien
trésaillir en aucun de mes élèves, ni même chez aucun des trois ou
quatre collègues de haut vol auxquels j'ai eu l'occasion d'en parler
de façon circonstanciée, que je fais un premier voyage de prospec-%
tion de ce ``monde nouveau'', de janvier à juin 1981. Ce premier jet se
matérialise en un paquet de quelques 1 300 pages manuscrites, bap-%
tisées ``La Longue Marche à travers la théorie de Galois''. Il s'agit
avant tout d'un effort de compréhension des relations entre groupes
de Galois ``arithmétiques'' et groupes fondamentaux profinis ``géomé-%
triques''. Assez vite il s'oriente vers un travail de formulation
calculatoire de l'opération de $\mathrm{Gal}(\bar{\mathbb{Q}}/\mathbb{Q})$ sur $\hat{\pi}_{0,3}$, et dans un
stade ultérieur, sur le groupe légèrement plus gros $S\ell(2,\mathbb{Z})^{\wedge}$, qui
donne lieu à un formalisme plus élégant et plus efficace. C'est au
cours de ce travail aussi (mais développé dans des notes distinctes)

\margin{93}

\clearpage
% Esquisse d'un programme (cote 119, copy 2): archive cote 119, site page 94 (PDF page 95), typescript page 16.
% typescript. Draft compared: esquisse/pages/esq_119_018.tex (copy 1, site 18).
% AI-transcribed by Claude Opus 5.5 (claude-opus-5-5, adaptive thinking, effort high), Message Batches API, batch esq119_full (msgbatch_015xqDHkTumDysRF1qngZhRJ), 2026-11-13; not yet checked by eye.
\scanpage{119/94}
\mspage{16}

qu'apparaît le thème central de la géométrie algébrique anabélienne,
qui est de reconstituer certaines variétés $X$ dites ``anabéliennes''
sur un corps absolu $K$ à partir de leur groupe fondamental mixte,
extension de $\mathrm{Gal}(\bar{K}/K)$ par $\pi_1(X_{\bar{K}})$ ; c'est alors que se dégage la
``conjecture fondamentale de la géométrie algébrique anabélienne'',
proche des conjectures de Mordell et de Tate que vient de démontrer
Faltings($^3$). C'est là aussi que s'amorcent une première réflexion sur
les groupes de Teichmüller, et les premières intuitions sur la struc-%
ture multiple de la ``tour de Teichmüller'' - les multiplicités modu-%
laires ouvertes $M_{g,\nu}$ apparaissant par ailleurs comme les premiers
exemples importants, en dimension $> 1$, de variétés (ou plutôt, de
multiplicités) qui semblent bien mériter l'appellation ``anabélienne''.
Vers la fin de cette période de réflexion, celle-ci m'apparaît com-%
me une réflexion fondamentale sur une théorie alors encore dans les
limbes, pour laquelle l'appellation ``\uline{théorie de Galois-Teichmüller}''
me semble plus appropriée que ``Théorie de Galois'' que j'avais d'abord
donnée à mes notes.

Ce n'est pas le lieu ici de donner un aperçu plus circons-%
tancié de cet ensemble de questions, intuitions, idées - y compris
des résultats palpables, certes. Le plus important me semble celui
signalé en passant au par. 2, savoir la fidélité de l'action exté-%
rieure de $\mathrm{I\Gamma} = \mathrm{Gal}(\bar{\mathbb{Q}}/\mathbb{Q})$ (et de ses sous-groupes ouverts) sur $\hat{\pi}_{0,3}$,
et plus généralement (si je me rappelle bien) sur le groupe fonda-%
mental de toute courbe algébrique ``anabélienne'' (i.e. dont le genre
$g$ et le ``nombre de trous'' $\nu$ satisfont l'inégalité $2g + \nu \geq 3$,
i.e. telle que $\chi(X) < 0$) définie sur une extension finie de $\mathbb{Q}$. Ce
résultat peut être considéré comme essentiellement équivalent au
théorème de Bielyi - c'est la première manifestation concrète, par
un énoncé mathématique précis, du ``message'' dont il a été question
plus haut.

Je voudrais terminer cet aperçu rapide par quelques mots
de commentaire sur la richesse vraiment inimaginable d'un groupe
anabélien typique comme le groupe $S\ell(2,\mathbb{Z})$ - sans doute le groupe
discret infini le plus remarquable qu'on ait rencontré, qui apparaît

\margin{94}

\clearpage
% Esquisse d'un programme (cote 119, copy 2): archive cote 119, site page 95 (PDF page 96), typescript page 17.
% typescript. Draft compared: esquisse/pages/esq_119_019.tex (copy 1, site 19).
% AI-transcribed by Claude Opus 5.5 (claude-opus-5-5, adaptive thinking, effort high), Message Batches API, batch esq119_full (msgbatch_015xqDHkTumDysRF1qngZhRJ), 2026-11-13; not yet checked by eye.
\scanpage{119/95}
\mspage{17}

sous une multiplicité d'avatars (dont certains ont été effleurés
dans le présent rapport), et qui du point de vue de la théorie de
Galois-Teichmüller peut être considéré comme la ``pierre de cons-%
truction'' fondamentale de la ``tour de Teichmüller''. L'élément de
structure de $S\ell(2,\mathbb{Z})$ qui me fascine avant tout, est bien sûr
l'action extérieure du groupe de Galois $\mathrm{I}\!\Gamma$ sur le compactifié
profini. Par le théorème de Bielyi, prenant les compactifiés pro-%
finis de sous-groupes d'indice fini de $S\ell(2,\mathbb{Z})$, et l'action ex-%
térieure induite (quitte à passer également à un sous-groupe ou-%
vert de $\mathrm{I}\!\Gamma$), \uline{on trouve essentiellement les groupes fondamentaux}
\uline{de toutes les courbes algébriques} (pas nécessairement compactes)
\uline{définis sur des corps de nombres} $K$, \uline{et l'action extérieure de}
$\mathrm{Gal}(\bar{K}/K)$ \uline{dessus} - du moins est-il vrai que tout tel groupe fon-%
\margin{?} damental apparaît comme quotient d'un des premiers groupes (*).
% NOTE: the typed (*) at the end of this line is circled by hand
Tenant compte du ``yoga anabélien'' (qui reste conjectural), disant
qu'une courbe algébrique anabélienne sur un corps de nombres $K$
(extension finie de $\mathbb{Q}$) est connu à isomorphisme près quand on
connaît son groupe fondamental mixte (ou ce qui revient au même,
l'action extérieure de $\mathrm{Gal}(\bar{K}/K)$ sur son groupe fondamental profini
géométrique), on peut donc dire que \uline{toutes les courbes algébriques}
\uline{définies sur des corps de nombres sont} ``\uline{contenues}'' \uline{dans le compac-}%
\uline{tifié profini} $S\ell(2,\mathbb{Z})\hat{}$, \uline{et dans la connaissance d'un certain sous-}%
\uline{groupe} $\mathrm{I}\!\Gamma$ \uline{du groupe des automorphismes extérieurs de ce dernier} !
Passant aux abélianisés des groupes fondamentaux précédents, on voit
notamment que toutes les représentations abéliennes $\ell$-adiques chères
à Tate et consorts, définies par des jacobiennes et jacobiennes gé-%
néralisées de courbes algébriques définies sur des corps de nombres,
sont contenues dans cette seule action de $\mathrm{I}\!\Gamma$ sur le groupe profini
anabélien $S\ell(2,\mathbb{Z})\hat{}$!($^{4}$)

\rule{5cm}{0.4pt}

(*) En fait, il s'agit de quotients de nature particulièrement tri-%
\margin{?} viale, par des sous-groupes abéliens produits de ``modules de Tate''
$\hat{\mathbb{Z}}(1)$, correspondants à des ``groupes-lacets'' autour de points à l'infini.

\margin{95}

\clearpage
% Esquisse d'un programme (cote 119, copy 2): archive cote 119, site page 96 (PDF page 97), typescript page 18.
% typescript. Draft compared: esquisse/pages/esq_119_020.tex (copy 1, site 20).
% AI-transcribed by Claude Opus 5.5 (claude-opus-5-5, adaptive thinking, effort high), Message Batches API, batch esq119_full (msgbatch_015xqDHkTumDysRF1qngZhRJ), 2026-11-13; not yet checked by eye.
\scanpage{119/96}
\mspage{18}

Il en est qui, face à cela, se contentent de hausser les
épaules d'un air  désabusé et de parier qu'il n'y a rien à tirer
de tout cela, sauf des rêves. Ils oublient, ou ignorent, que notre
science, et toute science, serait bien peu  de choses, si depuis
ses origines elle n'avait été nourrie des rêves et des visions
de ceux qui s'y adonnent avec passion.

4.        Dès le début de ma réflexion sur les cartes bidimension-%
nelles, je me suis intéressé plus particulièrement aux cartes dites
``\uline{régulières}'', c'est-à-dire celles dont le groupe des automorphismes
opère transitivement (et de ce fait, de façon simplement transitive)
sur l'ensemble des repères. Dans le cas orienté et en termes de
l'interprétation algébrico-géométrique du paragraphe précédent,
ce sont les cartes qui correspondent à un revêtement \uline{galoisien} de
la droite projective. Très vite aussi, et dès avant même qu'appa-%
raisse le lien avec la géométrie algébrique, il apparaît néces-%
saire aussi de ne pas exclure les cartes infinies, qui intervien-%
nent notamment de façon naturelle comme revêtements universels des
cartes finies. Il apparaît (comme conséquence immédiate du ``diction-%
naire'' des cartes, étendu au cas des cartes pas nécessairement finies)
que pour tout couple d'entiers naturels  $p, q$ \hand{$\geq$} 1, il existe à iso-%
morphisme (non unique) près une carte 1-\uline{connexe} et une seule qui
soit de type  $(p,q)$  i.e. dont tous les sommets soient d'ordre  $p$
et toutes les faces d'ordre  $q$, et cette carte est une carte régu-%
lière. Elle se trouve épinglée par le choix d'un repère, et son
groupe des automorphismes est alors canoniquement isomorphe au quo-%
tient du groupe cartographique (resp. du groupe cartographique
orienté, dans le cas orienté) par les relations supplémentaires
\[ \rho_s^{\,p} = \rho_f^{\,q} = 1. \]
Le cas où ce groupe est fini est le cas ``pythagoricien'' des cartes
régulières sphériques, le cas où il est infini donne les pavages

\margin{96}

\clearpage
% Esquisse d'un programme (cote 119, copy 2): archive cote 119, site page 97 (PDF page 98), typescript page 19.
% typescript. Draft compared: esquisse/pages/esq_119_021.tex (copy 1, site 21).
% AI-transcribed by Claude Opus 5.5 (claude-opus-5-5, adaptive thinking, effort high), Message Batches API, batch esq119_full (msgbatch_015xqDHkTumDysRF1qngZhRJ), 2026-11-13; not yet checked by eye.
\scanpage{119/97}
\mspage{19}

réguliers du plan euclidien ou du plan hyperbolique (*). Le lien
de la théorie combinatoire avec la théorie ``conforme'' des pavages
réguliers du plan hyperbolique était pressenti, avant qu'appa-%
raisse celui des cartes finies avec les revêtements finis de la
droite projective. Une fois ce lien compris, il devient évident
qu'il doit s'étendre également aux cartes infinies (régulières
ou non) : \uline{toute carte finie ou non, se réalise canoniquement sur
une surface conforme} (compacte si et seulement si la carte est
finie), \uline{en tant que revêtement ramifié de la droite projective com-%
plexe, ramifié seulement en les points} 0, 1, $\infty$. La seule difficulté
ici était de mettre au point le dictionnaire entre cartes topolo-%
giques et ensembles à opérateurs, qui posait quelques problèmes
conceptuels dans le cas infini, à commencer par la notion même de
``carte topologique''. Il apparaît nécessaire notamment, tant par
raison de cohérence interne du dictionnaire, que pour ne pas laisser
échapper certains cas intéressants de cartes infinies, de ne pas
exclure des sommets et des faces d'ordre infini. Ce travail de fon-%
dements a été fait également par J. Malgoire et C. Voisin, sur la
lancée de leur premier travail sur les cartes finies, et leur théo-%
rie fournit en effet tout ce qu'on était en droit d'en attendre
(et même plus...).

C'est en 1977 et 1978, parallèlement à deux cours de C 4
sur la géométrie du cube et sur celle de l'icosaèdre, que j'ai com-%
mencé à m'intéresser aux polyèdres réguliers, qui m'apparaissent
alors comme des ``réalisations géométriques'' particulièrement con-%
crètes de cartes combinatoires, les sommets, arêtes et faces étant
réalisés respectivement comme des points, des droites et des plans
dans un espace affine tridimensionnel convenable, avec respect des
relations d'incidence. Cette notion de réalisation géométrique

\rule{4cm}{0.4pt}

(*) Dans ces énoncés, il y a lieu de ne pas exclure le cas où $p$, $q$
peuvent prendre la valeur $+\infty$, qu'on rencontre notamment de façon
très naturelle comme pavages associés à certains polyèdres réguliers
infinis, cf. plus bas.

\margin{97}

\clearpage
% Esquisse d'un programme (cote 119, copy 2): archive cote 119, site page 98 (PDF page 99), typescript page 20.
% typescript. Draft compared: esquisse/pages/esq_119_022.tex (copy 1, site 22).
% AI-transcribed by Claude Opus 5.5 (claude-opus-5-5, adaptive thinking, effort high), Message Batches API, batch esq119_full (msgbatch_015xqDHkTumDysRF1qngZhRJ), 2026-11-13; not yet checked by eye.
\scanpage{119/98}
\mspage{20}

d'une carte combinatoire garde un sens sur un corps de base, et
même sur un anneau de base arbitraire. Elle garde également un
sens pour des polyèdres réguliers de dimension quelconque, en
remplaçant le groupe cartographique $\underline{C}_2$ par une variante n-dimen-%
sionnelle $\underline{C}_n$ convenable. Le cas n = 1, i.e. la théorie des
polygônes réguliers en caractéristique quelconque, fait l'objet
d'un cours DEA en 1977/78, et fait apparaître déjà quelques phé-%
nomènes nouveaux, comme aussi l'utilité de travailler non pas
dans un espace ambiant affine (ici le plan affine), mais dans
un espace \uline{projectif}. Ceci est dû notamment au fait que dans cer-%
taines caractéristiques (et notamment en caractéristique 2) le
centre d'un polyèdre régulier est rejeté à l'infini. D'autre
part, le contexte projectif, contrairement au contexte affine,
permet de développer avec aisance un formalisme de dualité pour
les polyèdres réguliers, correspondant au formalisme de dualité
des cartes combinatoires ou topologiques (où le rôle des sommets
et des faces, dans le cas n = 2 disons, se trouve interchangé).
Il se trouve que pour tout polyèdre régulier projectif, on peut
définir un hyperplan canonique associé, qui joue le rôle d'un
hyperplan à l'infini canonique, et permet de considérer le polyè-%
dre donné comme un polyèdre régulier affine.

L'extension de la théorie des polyèdres réguliers (et
plus généralement, de toutes sortes de configurations géométrico-%
combinatoires, y compris les systèmes de racines...) du corps de
base $\mathbb{R}$ ou $\mathbb{C}$ vers un anneau de base général, me semble d'une
portée comparable, dans cette partie de la géométrie, à l'exten-%
sion analogue qui a eu lieu depuis le début du siècle en géométrie
algébrique, ou depuis une vingtaine d'années en topologie ( *),
avec l'introduction du langage des schémas et celui des topos.

\rule{4cm}{0.4pt}

(*) En écrivant cela, je suis conscient que rares sont les topo-%
logues, encore aujourd'hui, qui se rendent compte de cet élargisse-%
ment conceptuel et technique de la topologie, et des ressources
qu'elle offre.

\margin{98}

\clearpage
% Esquisse d'un programme (cote 119, copy 2): archive cote 119, site page 99 (PDF page 100), typescript page 21.
% typescript. Draft compared: esquisse/pages/esq_119_023.tex (copy 1, site 23).
% AI-transcribed by Claude Opus 5.5 (claude-opus-5-5, adaptive thinking, effort high), Message Batches API, batch esq119_full (msgbatch_015xqDHkTumDysRF1qngZhRJ), 2026-11-13; not yet checked by eye.
\scanpage{119/99}
\mspage{21}

Ma réflexion sporadique sur cette question, pendant quelques années,
s'est bornée à dégager quelques principes de base simples, en at-%
tachant d'abord mon attention au cas des polyèdres réguliers
\uline{épinglés}, ce qui réduit à un minimum le bagage conceptuel néces-%
saire, et élimine pratiquement les questions de rationalité tant
soit peu délicates. Pour un tel polyèdre, on trouve une base (ou
repère) canonique de l'espace affine ou projectif ambiant, de
telle façon que les opérations du groupe cartographique $\underline{C}_n$, en-%
gendré par les réflexions fondamentales $\sigma_i$ $(0 \leq i \leq n)$, s'y
écrivent par des formules universelles, en termes de $n$ paramètres
$\alpha_1,\ldots,\alpha_n$ , qui géométriquement s'interprètent comme les doubles
des cosinus des ``angles fondamentaux'' du polyèdre. Le polyèdre se
reconstitue à partir de cette action, et du drapeau affine ou
projectif associé à la base choisie, en transformant ce drapeau
par tous les éléments du groupe engendré par les réflexions fon-%
damentales. Ainsi le n-polyèdre épinglé ``universel'' est-il défini
canoniquement sur l'anneau de polynômes à $n$ indéterminées
\[ \mathbb{Z}[\underline{\alpha}_1,\ldots,\underline{\alpha}_n] , \]
ses spécialisations sur des corps de base arbitraires $k$ (via
des valeurs $\alpha_i \in k$ données aux indéterminées $\underline{\alpha}_i$) donnant des
polyèdres réguliers correspondant à des types combinatoires di-%
vers. Dans ce jeu, il n'est pas question de se borner à des po-%
lyèdres réguliers finis, ni même à des polyèdres réguliers dont
les facettes soient d'ordre fini, i.e. pour lesquels les para-%
mètres $\alpha_i$ soient des racines d'équations ``semicyclotomiques''
convenables, exprimant que les ``angles fondamentaux'' (dans le
cas où le corps de base est $\mathbb{R}$) sont commensurables à $2\pi$. Déjà
quand $n = 1$, le polygône régulier peut-être le plus intéressant
de tous (moralement celui du polygone régulier à un seul côté !)
est celui qui correspond à $\alpha = 2$, donnant lieu à une conique cir-%
conscrite parabolique, i.e. tangente à la droite à l'infini.

Le cas fini est celui où le groupe engendré par les réflexions
fondamentales, qui est aussi le groupe des automorphismes du

\margin{99}

\clearpage
% Esquisse d'un programme (cote 119, copy 2): archive cote 119, site page 100 (PDF page 101), typescript page 22.
% typescript. Draft compared: esquisse/pages/esq_119_024.tex (copy 1, site 24).
% AI-transcribed by Claude Opus 5.5 (claude-opus-5-5, adaptive thinking, effort high), Message Batches API, batch esq119_full (msgbatch_015xqDHkTumDysRF1qngZhRJ), 2026-11-13; not yet checked by eye.
\scanpage{119/100}
\mspage{22}

polyèdre régulier envisagé, est fini. Dans le cas du corps de
base $\mathbb{R}$ (ou $\mathbb{C}$, ce qui revient au même), et pour $n = 2$, les cas
finis sont bien connus depuis l'antiquité - ce qui n'exclut pas
que le point de vue schématique y fasse apparaître des charmes
nouveau ; on peut dire cependant qu'en spécialisant l'isosaèdre
(par exemple) sur des corps de base finis de caractéristique
arbitraire, c'est toujours un icosaèdre, avec sa combinatoire
propre et le même groupe d'automorphismes simple d'ordre 60 qu'on
obtient. La même remarque s'applique aux polyèdres réguliers finis
de dimension supérieure, étudiés de façon systématique dans deux
beaux livres de Coxeter. La situation est toute autre si on part
d'un polyèdre régulier \uline{infini}, sur un corps tel que $\mathbb{Q}$ disons,
et qu'on le ``spécialise'' sur les corps premiers $\mathbb{F}_p$ (opération
bien définie pour tout $p$ sauf un nombre fini de nombres premiers).
Il est clair que tout polyèdre régulier sur un corps fini est fini -
\uline{on trouve donc une infinité de polyèdres réguliers finis pour} $p$
\uline{variable}, \uline{dont le type combinatoire}, ou ce qui revient au même,
le groupe des automorphismes, \uline{varie de façon} ``\uline{arithmétique}'' avec p.
Cette situation est particulièrement intrigante  dans le cas où
$n = 2$, où on dispose de la relation explicitée au paragraphe pré-%
cédent entre 2-cartes combinatoires, et courbes algébriques dé-%
finies sur des corps de nombres. Dans ce cas, un polyèdre régu-%
lier infini défini sur un corps infini quelconque (et de ce fait,
sur une sous-$\mathbb{Z}$-algèbre à deux générateurs de celui-ci) donne donc
naissance à une infinité de courbes algébriques définies sur des
corps de nombres, qui sont des revêtements galoisiens ramifiés
seulement en 0, 1, $\infty$ de la droite projective standard. Le cas opti-%
mum est bien sûr celui où on part du 2-polyèdre régulier universel,
ou plutôt de celui qui s'en déduit par passage au corps des frac-%
tions $\mathbb{Q}(\alpha_1,\alpha_2)$ de son anneau de base. Ceci soulève une foule de
questions nouvelles, aussi bien des vagues que des précises, dont
je n'ai eu loisir encore d'examiner de plus près aucune - je ne
citerai que celle-ci : quelles sont exactement les 2-cartes régu-%
lières finies, ou ce qui revient au même, les groupes quotients

\margin{100}

\clearpage
% Esquisse d'un programme (cote 119, copy 2): archive cote 119, site page 101 (PDF page 102), typescript page 23.
% typescript. Draft compared: esquisse/pages/esq_119_025.tex (copy 1, site 25).
% AI-transcribed by Claude Opus 5.5 (claude-opus-5-5, adaptive thinking, effort high), Message Batches API, batch esq119_full (msgbatch_015xqDHkTumDysRF1qngZhRJ), 2026-11-13; not yet checked by eye.
\scanpage{119/101}
\mspage{23}

finis du groupe 2-cartographique, qui proviennent de 2-polyèdres
réguliers sur des corps finis (*) ? Les obtiendrait-on toutes,
et si oui : comment ?

Ces réflexions font apparaître en pleine lumière ce fait,
qui pour moi était entièrement inattendu, que la théorie des
polyèdres réguliers finis, déjà dans le cas de la dimension
$n = 2$, est infiniment plus riche, et notamment donne infiniment
plus de formes combinatoires différentes, dans le cas où on admet
des corps de base de caractéristique non nulle,\quad que dans le
cas considéré jusqu'à présent, où les corps de base étaient res-%
treint à $\mathbb{R}$ ,ou à la rigueur $\mathbb{C}$ (dans le cas de ce que Coxeter
appelle des ``polyèdres réguliers complexes'', et que je préfère
appeler ``pseudo-polyèdres réguliers définis sur $\mathbb{C}$ '') (**). De plus,
il semble que cet élargissement du point de vue doive aussi jeter
un jour nouveau sur les cas déjà connus. Ainsi, examinant l'un
après l'autre les polyèdres pythagoriciens, j'ai vu se répéter
à chaque fois un même petit miracle, que j'ai appelé le \uline{paradigme
combinatoire} du polyèdre envisagé. Vaguement parlant, il peut se
décrire en disant que lorsqu'on regarde la spécialisation du po-%
lyèdre dans la caractéristique , ou l'une des caractéristiques, la
(ou les) plus singulière(s) (ce sont les caractéristique 2 et 5

\rule{5cm}{0.4pt}

(*) Ce sont les mêmes d'ailleurs que ceux provenant de polyèdres
réguliers sur des corps quelconques, ou algébriquement clos, comme
on voit par des arguments de spécialisation standards.

(**) Les pseudo-polyèdres épinglés se décrivent de la même façon que
les polyèdres épinglés, avec cette seule différence que les réflexions
fondamentales $\sigma_i$ $(0 \leq i \leq n)$ sont remplacées ici par des \uline{pseudo-%
réflexions} (que Coxeter suppose de plus d'ordre fini, comme il se
borne aux structures combinatoires finies). Cela conduit simplement
à introduire pour chacun des $\sigma_i$ un invariant numérique supplémen-%
taire $\beta_i$ , de sorte que le n-pseudo-polyèdre universel peut se dé-%
finir encore sur un anneau de polynômes à coefficients entiers, en
les n+(n+1) variables $\underline{\alpha}_i$ $(1 \leq i \leq n)$ et $\underline{\beta}_j$ $(0 \leq j \leq n)$.

\margin{101}

\clearpage
% Esquisse d'un programme (cote 119, copy 2): archive cote 119, site page 102 (PDF page 103), typescript page 24.
% typescript. Draft compared: esquisse/pages/esq_119_026.tex (copy 1, site 26).
% AI-transcribed by Claude Opus 5.5 (claude-opus-5-5, adaptive thinking, effort high), Message Batches API, batch esq119_full (msgbatch_015xqDHkTumDysRF1qngZhRJ), 2026-11-13; not yet checked by eye.
\scanpage{119/102}
\mspage{24}

pour l'icosaèdre, la caractéristique 2 pour l'octaèdre), on lit,
sur le polyèdre régulier géométrique sur le corps fini concerné
($\mathbb{F}_4$ et $\mathbb{F}_5$ pour l'icosaèdre, $\mathbb{F}_2$ pour l'octaèdre) une descrip-%
tion particulièrement élégante (et inattendue) de la combina-%
toire du polyèdre. Il m'a semblé même entrevoir là un principe
d'une grande généralité, que j'ai crû retrouver notamment dans
une réflexion ultérieure sur la combinatoire du système des 27
droites d'une surface cubique, et ses relations avec le système
de racines $E_7$. Qu'un tel principe existe bel et bien et qu'on
réunisse même à le dégager de son manteau de brumes, ou qu'il
recule au fur et à mesure où on le poursuit et qu'il finisse par
s'évanouir comme une Fata Morgana, j'y trouve pour ma part une
force de motivation, une fascination peu communes, comme celle du
rêve peut-être. Nul doute que de suivre un tel appel de l'infor-%
mulé, de l'informe qui cherche forme, d'un entrevu
élusif qui semble prendre plaisir à la fois à se dérober et à se
manifester - ne peut que mener loin, alors que nul ne pourrait
prédire, où ...

Pourtant, pris par d'autres intérêts et tâches, je n'ai
pas jusqu'à présent suivi cet appel, ni rencontré personne d'autre
qui ait voulu l'entendre, et encore moins le suivre. Mis à part
quelques digressions vers d'autres types de structures géométrico-
combinatoires, mon travail ici encore s'est borné à un premier tra-%
vail de dégrossissage et d'intendance, sur lequel il est inutile
de m'étendre plus ici\hand{$^{(5)}$}. Le seul point qui peut-être mérite encore
mention, est l'existence et l'unicité de l'hyperquadrique circons-%
crite à un n-polyèdre régulier donné, dont l'équation peut s'expli-%
citer par des formules simples en termes des paramètres fondamentaux
$\alpha_i$ (*). Le cas qui m'intéresse le plus est celui où $n = 2$, et le

\rule{5cm}{0.4pt}

(*) Un résultat analogue vaut pour les pseudo-polyèdres. Il semblerait
que les ``caractéristiques exceptionnelles'' dont il a été question plus
haut, pour les spécialisations d'un polyèdre donné, sont celles pour
\margin{x} lesquelles l'hyperquadr\deleted{\unsure{at}}ique circonscrite est, soit dégénérée, soit
tangente à l'hyperplan à l'infini.

\margin{102}

\clearpage
% Esquisse d'un programme (cote 119, copy 2): archive cote 119, site page 103 (PDF page 104), typescript page 25.
% typescript. Draft compared: esquisse/pages/esq_119_027.tex (copy 1, site 27).
% AI-transcribed by Claude Opus 5.5 (claude-opus-5-5, adaptive thinking, effort high), Message Batches API, batch esq119_full (msgbatch_015xqDHkTumDysRF1qngZhRJ), 2026-11-13; not yet checked by eye.
\scanpage{119/103}
\mspage{25}

temps me semble mûr pour réécrire une version nouvelle, en style
moderne, du classique livre de Klein sur l'icosaèdre et les autres
polyèdres pythagoriciens. Ecrire un tel exposé sur les 2-polyèdres
réguliers serait une magnifique occasion pour un jeune chercheur
de se familiariser aussi bien avec la géométrie des polyèdres et
leurs liens avec les géométries sphèrique, euclidienne, hyperbo-%
lique, et avec les courbes algébriques, qu'avec le langage et les
techniques de base de la géométrie algébrique moderne. S'en trou-%
vera-t-il un un jour pour saisir cette occasion ?

5. \qquad Je voudrais maintenant dire quelques mots sur certaines
réflexions qui m'ont fait comprendre le besoin de fondements nou-%
veaux pour la topologie ``géométrique'', dans une direction toute
différente de la notion de topos, et indépendante même des besoins
de la géométrie algébrique dite ``abstraite'' (sur des corps et an-%
neaux de base généraux). Le problème de départ, qui a commencé à
m'intriguer il doit y avoir une quinzaine d'années déjà, était
celui de définir une théorie de ``dévissage'' des structures stra-%
tifiées, pour les reconstituer, par un procédé canonique, à partir
de ``pièces de construction'' canoniquement déduits de la structure
donnée. Probablement l'exemple principal qui m'avait alors amené
à cette question était celui de la stratification canonique d'une
variété algèbrique singulière (ou d'un espace analytique complexe
ou réel singulier) par la suite décroissante de ses ``lieux singuliers''
successifs. Mais je devais sans doute pressentir déjà l'ubiquité
des structures stratifiées dans pratiquement tous les domaines de
la géométrie (que d'autres sûrement ont vu clairement bien avant
moi). Depuis, j'ai vu apparaître de telles structures, notamment,
dans toute situation de ``modules'' pour des objets géométriques sus-%
ceptibles non seulement de variation continue, mais en même temps
de phénomènes de ``degénérescence'' (ou de ``spécialisation'') - les
strates correspondant alors aux divers ``niveaux de singularité''
(ou aux types combinatoires associés) pour les objets considérés.

\margin{103}

\clearpage
% Esquisse d'un programme (cote 119, copy 2): archive cote 119, site page 104 (PDF page 105), typescript page 26.
% typescript. Draft compared: esquisse/pages/esq_119_028.tex (copy 1, site 28).
% AI-transcribed by Claude Opus 5.5 (claude-opus-5-5, adaptive thinking, effort high), Message Batches API, batch esq119_full (msgbatch_015xqDHkTumDysRF1qngZhRJ), 2026-11-13; not yet checked by eye.
\scanpage{119/104}
\mspage{26}

Les multiplicités modulaires compactifiées $\hat{M}_{g,\nu}$ de Mumford-%
Deligne pour les courbes algébriques stables de type $(g,\nu)$ en
fournissent un exemple typique et particulièrement inspirant, qui
a joué un rôle de motivation important dans la reprise de ma ré-%
flexion sur les structures stratifiées , de décembre 1981 à janvier
1982. La géométrie bidimensionelle fournit de nombreux autres
exemples de telles structures stratifiées modulaires, qui toutes
d'ailleurs (sauf expédients de \unsure{rigidifidation|rigidification}), apparaissent
comme des ``multiplicités'' plutôt que comme des espaces ou variétés
au sens ordinaire (les points de ces multiplicité pouvant avoir
des groupes d'automorphismes non triviaux). Parmi les objets de
géométrie bidimensionelle donnant lieu à de telles structures mo-%
dulaires stratifiées de dimension arbitraire, voire de dimension
infinie, je citerai les polygônes (euclidiens, ou sphériques, ou
hyperboliques), les systèmes de droites dans un plan (projectif
disons), les systèmes de ``pseudodroites'' dans un plan projectif
topologique, ou les courbes immergées à \unsure{crc isements|croisements} normaux plus
générales, dans une surface (compacte disons) donnée.

L'exemple non trivial le plus simple d'une structure
stratifiée s'obtient en considérant une paire $(X,Y)$ d'un espace
$X$ et d'un sous-espace fermé $Y$, en faisant une hypothèse d'équi-%
singularité convenable de $X$ le long de $Y$, et en supposant de
plus (pour fixer les idées) que les deux strates $Y$ et $X\setminus Y$ sont
des \uline{variétés} topologiques. L'idée naïve, dans une telle situation,
est de prendre ``le'' voisinage tubulaire $T$ de $Y$ dans $X$, dont
le bord $\partial T$ devrait être une variété lisse également, fibrée à
fibres lisses et compactes sur $Y$, $T$ lui-même s'identifi\hand{\unsure{/}}ant au
fibré en cônes sur $T$ associé au fibré précédent. Posant
\[ U = X \setminus \mathrm{Int}(T), \]
on trouve une variété à bord dont le bord est canoniquement isomor-%
phe à celui de $T$. Ceci dit, les ``pièces de construction'' prévues
sont la variété à bord $U$ (compacte si $X$ était compact, et qui
remplace en la précisant la strate ``ouverte'' $X\setminus Y$) et la variété

\margin{104}

\clearpage
% Esquisse d'un programme (cote 119, copy 2): archive cote 119, site page 105 (PDF page 106), typescript page 27.
% typescript. Draft compared: esquisse/pages/esq_119_029.tex (copy 1, site 29).
% AI-transcribed by Claude Opus 5.5 (claude-opus-5-5, adaptive thinking, effort high), Message Batches API, batch esq119_full (msgbatch_015xqDHkTumDysRF1qngZhRJ), 2026-11-13; not yet checked by eye.
\scanpage{119/105}
\mspage{27}

(sans bord)  Y, avec comme structure supplémentaire les reliant
l'application dite ``de recollement''
\[ f : \partial U \longrightarrow Y \]
qui est une fibration propre et lisse. La situation de départ
(X,Y)  se reconstitue à partir de  $(U,Y, f : \partial U \longrightarrow Y)$  par la
formule
\[ X \simeq U \amalg_{\partial U} Y \]
(somme amalgamée sous  $\partial U$, s'envoyant dans  U  et  Y  via l'inclu-%
sion resp. l'application de recollement).

Cette vision naïve se heurte immédiatement à des dif-%
ficultés diverses. La première est la nature un peu vague de la
notion même de voisinage tubulaire, qui ne prend un sens tant soit
peu précis qu'en présence de structures beaucoup plus rigides que
la seule structure topologique, telles la structure ``linéaire par
morceaux'', ou riemanienne (plus généralement, d'espace avec fonc-%
tion distance) ; l'ennui ici est que dans aucun des exemples
auxquels on pense spontanément, on ne dispose naturellement d'une
structure de ce type - tout au mieux d'une classe d'équivalence
de telles structures, permettant de rigidifier un tantinet la
situation. Si par ailleurs on admet qu'on a pu trouver un expé-%
dient pour trouver un voisinage tubulaire ayant les propriétés
voulues, qui de plus soit unique modulo un automorphisme (topolo-%
gique, disons) de la situation, automorphisme qui de plus res-%
pecte la structure fibrée fournie par la fonction de recollement,
il reste la difficulté de la non-canonicité des choix faits, l'au-%
tomorphisme en question n'étant visiblement pas unique, quoi qu'on
fasse pour le ``normaliser''. L'idée ici, pour rendre canonique ce
qui ne l'est pas, est de travailler systématiquement dans des
``\uline{catégories isotopiques}'' associées aux catégories de nature topolo-%
gique s'introduisant dans ces questions (telle la catégorie des
paires admissibles  (X,Y)  et des homéomorphisme de telles paires,
etc), en gardant les mêmes objets, mais en prenant comme ``morphismes''

\margin{105}

\clearpage
% Esquisse d'un programme (cote 119, copy 2): archive cote 119, site page 106 (PDF page 107), typescript page 28.
% typescript. Draft compared: esquisse/pages/esq_119_030.tex (copy 1, site 30).
% AI-transcribed by Claude Opus 5.5 (claude-opus-5-5, adaptive thinking, effort high), Message Batches API, batch esq119_full (msgbatch_015xqDHkTumDysRF1qngZhRJ), 2026-11-13; not yet checked by eye.
\scanpage{119/106}
\mspage{28}

les classes d'isotopie (dans un sens dicté sans ambiguité par le
contexte) d'isomorphismes (voire même, de morphismes plus généraux
que des isomorphismes). Cette idée, qui est reprise avec succès
dans la thèse de Yves Ladegaillerie notamment (cf. début du par. 3),
m'a servi de façon systématique dans toutes mes réflexions ulté-%
rieures de topologie combinatoire, quand il s'est agi de formuler
avec précision des théorèmes de traduction de situations topolo-%
giques, en termes de situations combinatoires. Dans la situation
actuelle, mon espoir était d'arriver à formuler (et à prouver !)
un théorème d'équivalence entre deux catégories isotopiques conve-%
nables, l'une étant la catégorie des ``paires admissibles'' $(X,Y)$,
l'autre celle des ``triples admissibles'' $(U,Y,f)$, où $Y$ est une
variété, $U$ une variété à bord, et $f : \partial U \longrightarrow Y$ une fibration
propre et lisse. De plus, bien sûr, j'espérais qu'un tel énoncé,
modulo un travail de nature essentiellement algébrique, s'éten-%
drait de lui-même en un énoncé plus sophistiqué, s'appliquant aux
structures stratifiées générales.

Très vite il apparaissait qu'il ne pouvait être ques-%
tion d'obtenir un énoncé aussi ambitieux dans le contexte des
espaces topologiques, à cause des sempiternels phénomènes de
``sauvagerie''. Déjà quand $X$ lui-même est une variété et $Y$ ré-%
duit à un point, on se bute à la difficulté que le cône sur un
espace compact $Z$ peut être une variété en son sommet, sans que
$Z$ soit homéomorphe à une sphère, ni même soit une variété. Il
était clair également que les contextes de structures plus rigides
qui existaient à l'époque, tel le contexte ``linéaire par morceaux'',
étaient également inadéquats - une des raisons rédhibitoires com-%
munes étant qu'ils ne permettaient pas, pour une paire $(U,S)$
d'un ``espace'' $U$ et d'un sous-espace fermé $S$, et une application
de recollement $f : S \longrightarrow T$, de construire la somme amalgamée cor-%
respondante. C'est quelques années plus tard que j'étais informé
de la théorie de Hironaka des ensembles qu'il appelle, je crois,
``semi-analytiques'' (réels), qui satisfont à certaines des condi-%
tions de stabilité essentielles (sans doute même à toutes) néces-

\margin{106}

\clearpage
% Esquisse d'un programme (cote 119, copy 2): archive cote 119, site page 107 (PDF page 108), typescript page 29.
% typescript. Draft compared: esquisse/pages/esq_119_031.tex (copy 1, site 31).
% AI-transcribed by Claude Opus 5.5 (claude-opus-5-5, adaptive thinking, effort high), Message Batches API, batch esq119_full (msgbatch_015xqDHkTumDysRF1qngZhRJ), 2026-11-13; not yet checked by eye.
\scanpage{119/107}
\mspage{29}

saires au développement d'un contexte utilisable de ``topologie
modérée''. Du coup cela relance une réflexion sur les fondements
d'une telle topologie, dont le besoin m'apparaît de plus en plus
clairement.

Avec un recul d'une dizaine d'années, je dirais aujour-%
d'hui, à ce sujet, que la ``\uline{topologie générale}'' \uline{a été développée}
(dans les années trente et quarante) \uline{par des analystes et pour}
\uline{les besoins de l'analyse}, non pour les besoins de la topologie
proprement dite, c'est à dire l'étude des \uline{propriétés topologiques}
\uline{de formes géométriques} diverses. Ce caractère inadéquat des fon-%
dements de la topologie se manifeste dès les débuts, par des
``faux problèmes'' (au point de vue au moins de l'intuition topo-%
logique des formes) comme celle de ``l'invariance du domaine'',
alors même que la solution de ce dernier par Brouwer l'amène à
introduire des idées géométriques nouvelles importantes. Aujour-%
d'hui encore, comme aux temps héroïques où on voyait pour la pre-%
mière fois et avec inquiétude des courbes remplir allègrement
des carrés et des cubes, quand on se propose de faire de la géo-%
métrie topologique dans le contexte technique des espaces topo-%
logiques, on se heurte à chaque pas à des difficultés parasites
tenant aux phénomènes sauvages. Ainsi, en dehors de cas de (très)
basse dimension, il ne peut guère être possible, pour un espace
donné $X$ (une variété compacte disons), d'étudier le type d'homo-%
topie (disons) du groupe des automorphismes de $X$, ou de l'espace
des plongements, ou immersions etc. de $X$ dans quelque autre
espace $Y$ - alors qu'on sent que ces invariants devraient faire
partie de l'arsenal des invariants essentiels associés à $X$, ou
au couple $(X,Y)$, etc, au même titre que l'espace fonctionnel
$\underline{\mathrm{Hom}}\,(X,Y)$ familier en topologie homotopique. Les topologues élu-%
dent la difficulté, sans l'affronter, en se rabattant sur des con-%
textes voisins du contexte topologique et moins marqués de sauva-%
gerie que lui, comme les variétés différentiables, les espaces PL
(linéaires par morceaux), etc, dont visiblement aucun n'est ``bon'',
i.e. n'est stable par les opérations topologiques les plus évidentes,

\margin{107}

\clearpage
% Esquisse d'un programme (cote 119, copy 2): archive cote 119, site page 108 (PDF page 109), typescript page 30.
% typescript. Draft compared: esquisse/pages/esq_119_032.tex (copy 1, site 32).
% AI-transcribed by Claude Opus 5.5 (claude-opus-5-5, adaptive thinking, effort high), Message Batches API, batch esq119_full (msgbatch_015xqDHkTumDysRF1qngZhRJ), 2026-11-13; not yet checked by eye.
\scanpage{119/108}
\mspage{30}

telles les opérations de contraction-recollement (sans même passer
à des opérations du type $X \longrightarrow \mathrm{Aut}(X)$ qui font quitter le para-%
dis des ``espaces'' de dimension finie). C'est là une façon de tour-%
ner autour du pot ! Cette situation, comme tant de fois déjà dans
l'histoire de notre science, met simplement en évidence cette iner-%
tie quasi-insurmontable de l'esprit, alourdi par des conditionne-%
ments d'un poids considérable, pour porter un regard sur une ques-%
tion de fondements,donc sur le contexte même dans lequel on vit,
respire, travaille - plutôt que de l'accepter comme un donné immu-%
able. C'est à cause de cette inertie sûrement qu'il a fallu des
millénaires pour qu'une idée ou une réalité aussi enfantine que
le zéro, un groupe, ou une forme topologique, trouve droit de cité
en mathématiques. C'est par elle aussi, sûrement, que le carcan
de la topologie générale continue à être traîné patiemment par
des générations de topologues, la ``sauvagerie'' étant portée comme
une fatalité inéluctable qui serait enracinée dans la nature même
des choses.

Mon approche vers des fondements possibles d'une topo-%
logie modérée a été une approche axiomatique. Plutôt que de dé-%
clarer (chose qui serait parfaitement raisonnable certes) que
les ``espaces modérés'' cherchés ne sont autres (disons) que les
espaces semianalytique de Hironaka, et de développer dès lors dans
ce contexte l'arsenal des constructions et notions familières en
topologie, plus celles certes qui jusqu'à présent n'avaient pu
être développés et pour cause, j'ai préféré m'attacher à dégager
ce qui, parmi les propriétés géométriques de la notion d'ensemble
semianalytique dans un espace $\mathbb{R}^{n}$, permet d'utiliser ceux-ci comme
``modèles'' locaux d'une notion ``\uline{d'espace modéré}'' (en l'occurence,
semianalytique), et ce qui (on l'espère !) rend cette notion d'es-%
pace modéré suffisamment souple, pour pouvoir bel et bien servir
de notion de base pour une ``topologie modérée'' propre à exprimer
avec aisance l'intuition topologique des formes. Ainsi, une fois
le travail de fondements qui s'impose accompli, il apparaîtra non
\uline{une} ``théorie modérée'', mais une vaste infinité, allant de la plus

\margin{108}

\clearpage
% Esquisse d'un programme (cote 119, copy 2): archive cote 119, site page 109 (PDF page 110), typescript page 31.
% typescript. Draft compared: esquisse/pages/esq_119_033.tex (copy 1, site 33).
% AI-transcribed by Claude Opus 5.5 (claude-opus-5-5, adaptive thinking, effort high), Message Batches API, batch esq119_full (msgbatch_015xqDHkTumDysRF1qngZhRJ), 2026-11-13; not yet checked by eye.
\scanpage{119/109}
\mspage{31}

stricte de toutes, celles des ``espaces $\bar{\mathbb{Q}}_r$-algébriques par morceaux''
(où $\bar{\mathbb{Q}}_r = \bar{\mathbb{Q}} \cap \mathbb{R}$), vers celle qui (à tort ou à raison) m'apparaît
comme probablement la plus vaste, savoir celle des ``espaces ana-%
litiques réels par morceaux'' (ou semianalytiques dans la termi-%
nologie de Hironaka). Parmi les théorèmes de fondements envisagés
dans mon programme, il y a un \uline{théorème de comparaison} qui, vague-%
ment parlant, dira qu'\uline{on trouvera essentiellement les mêmes caté-%
gories isotopiques} (ou même $\infty$-isotopiques), quelle que soit la
théorie modérée avec laquelle on travaille\hand{$^{(6)}$}. De façon plus précise,
il s'agit de mettre le doigt sur un système d'axiomes suffisam-%
ment riche, pour impliquer (entre bien autres choses !) que si on
a deux théories modérées \hand{$\mathcal{T}$}, \hand{$\mathcal{T}$}' avec \hand{$\mathcal{T}$} plus fine que \hand{$\mathcal{T}$}' (dans
un sens évident), et si X, Y sont deux espaces T-modérés, qui
définissent donc aussi des espaces T'-modérés correspondants,
l'application canonique
\[ \underline{\mathrm{Isom}}_{\hand{\mathcal{T}}}(X,Y) \longrightarrow \underline{\mathrm{Isom}}_{\hand{\mathcal{T}}'}(X,Y) \]
induit une bijection sur l'ensemble des composantes connexes (ce
qui impliquera que la catégorie isotopique des T-espaces est équi-%
valente à celle des T'-espaces), et même, est une équivalence
d'homotopie (ce qui signifie qu'on a même une équivalence pour
les catégories ``$\infty$-isotopique'', plus fines que les catégories
isotopiques où on ne retient que le $\pi_0$ des espaces d'isomorphis-%
mes). Ici les \uline{Isom} peuvent être définis de façon évidente comme
ensembles semisimpliciaux par exemple, pour pouvoir donner un sens
précis à l'énoncé précédent. Des énoncés analogues devraient être
vrais, en remplaçant les ``espaces'' \uline{Isom} par d'autres espaces d'ap-%
plications, soumises à des conditions géométriques standard, comme
celle d'être des plongements, des immersions, lisses, étales, des
fibrations etc. Egalement, on s'attend à avoir des énoncés analo-%
gues, où X, Y sont remplacés par des systèmes d'espaces modérés,
tels ceux qui interviennent dans une théorie de dévissage des
structures stratifiées - de telle sorte que dans un sens techni-%
que précis, cette théorie de dévissage sera, elle aussi, essen-%

\margin{109}

\clearpage
% Esquisse d'un programme (cote 119, copy 2): archive cote 119, site page 110 (PDF page 111), typescript page 32.
% typescript. Draft compared: esquisse/pages/esq_119_034.tex (copy 1, site 34).
% AI-transcribed by Claude Opus 5.5 (claude-opus-5-5, adaptive thinking, effort high), Message Batches API, batch esq119_full (msgbatch_015xqDHkTumDysRF1qngZhRJ), 2026-11-13; not yet checked by eye.
\scanpage{119/110}
\mspage{32}

tiellement indépendante de la théorie modérée choisie pour l'ex-%
primer.

Le premier test décisif pour un bon système d'axiomes
sur une notion de ``partie modérée de $\mathbb{R}^n$ '' me semble la possibi-%
lité de prouver de tels théorèmes de comparaison. Je me suis con-%
tenté jusqu'à présent de dégager un système d'axiomes plausible
provisoire, sans avoir aucune assurance qu'il ne faudra y rajou-%
ter d'autres axiomes, que seul un ``travail sur pièces'' sans doute
permettra de faire apparaître. Le plus fort des axiomes que j'ai
introduits, et celui sans doute dont la vérification dans les cas
d'espèce est (ou sera) la plus délicate, est un \uline{axiome de trian-}%
\uline{gulabilité} (modérée, il va sans dire) d'une partie modérée de $\mathbb{R}^n$.
Je ne me suis pas essayé à prouver en termes de ces seuls axiomes
le théorème de comparaison, j'ai eu l'impression néanmoins (à tort
ou à raison encore !) que cette démonstration, qu'elle nécessite
ou non l'introduction de quelque axiome supplémentaire, ne pré-%
sentera pas de grosse difficulté technique. Il est bien possible
que les difficultés au niveau technique, pour le développement
de fondements satisfaisants de la topologie modérée, y inclus une
théorie de dévissage des structures modérées stratifiées, soient
déjà pour l'essentiel concentrées dans les axiomes, et par suite
essentiellement surmontées dès à présent par des théorèmes de
triangulabilité à la Lojasiewicz et Hironaka. Ce qui fait défaut,
encore une fois, n'est nullement la virtuosité technique des ma-%
thématiciens, parfois impressionnante, mais l'audace (ou simple-%
mement l'innocence ...) pour s'affranchir d'un contexte familier
accepté par un consensus sans failles ...

Les avantages d'une approche axiomatique vers des fon-%
dements de la topologie modérée me semblent assez évidents. Ainsi,
pour considérer une variété algébrique complexe, ou l'ensemble des
points réels d'une variété algébrique définie sur $\mathbb{R}$, comme un
espace modéré, il semble préférable de travailler dans la théorie
``$\mathbb{R}$-algébrique par morceaux'', voire même dans la théorie $\overline{\mathbb{Q}}_r$-algébrique
par morceaux (où $\overline{\mathbb{Q}}_r = \overline{\mathbb{Q}} \cap \mathbb{R}$) quand il s'agit de variétés définies

\margin{110}

\clearpage
% Esquisse d'un programme (cote 119, copy 2): archive cote 119, site page 111 (PDF page 112), typescript page 33.
% typescript. Draft compared: esquisse/pages/esq_119_035.tex (copy 1, site 35).
% AI-transcribed by Claude Opus 5.5 (claude-opus-5-5, adaptive thinking, effort high), Message Batches API, batch esq119_full (msgbatch_015xqDHkTumDysRF1qngZhRJ), 2026-11-13; not yet checked by eye.
\scanpage{119/111}
\mspage{33}

sur des corps de nombres, etc. L'introduction d'un sous-corps
$K \subset \mathbb{R}$ associé à la théorie $\hand{\mathcal{T}}$ (formé des points de $\mathbb{R}$ qui sont
$T$-\uline{modérés}, i.e. tels que l'ensemble uniponctuel correspondant le
soit) permet d'introduire pour tout point $x$ d'un espace modéré
$X$ un corps résidu\deleted{a}\inserted{e}l $k(x)$, qui est une sous-extension de $\mathbb{R}/K$
algébriquement fermée dans $\mathbb{R}$, et de degré de transcendance fini
sur $K$ (majoré par la dimension topologique de $X$). Quand le
degré de transcendance de $\mathbb{R}$ sur $K$ est infini, on trouve une
notion de degré de transcendance (ou ``dimension'') d'un point d'un
espace modéré, voisin de la notion familière en géométrie algé-%
brique. De telles notions sont absentes dans la topologie modé-%
rée ``semianalytique'', qui par contre apparaît comme le contexte
topologique tout indiqué pour inclure les espaces analytiques réels
et complexes.

Parmi les premiers théorèmes auxquels on s'attend dans
\margin{/}une topologie modérée comme je l'entrevois, mais\hand{/} à part les théo-%
rèmes de comparaison, sont les énoncés qui établissent, dans un
sens convenable, l'existence et l'unicité ``du'' voisinage tubulaire
d'un sous-espace modéré fermé dans un espace modéré (compact pour
simplifier), les façons concrètes de l'obtenir (par exemple à par-%
tir de toute application modéré $X \longrightarrow \mathbb{R}^{+}$ admettant $Y$ comme
ensemble de ses zéros), la description de son ``bord'' (alors qu'en
général ce n'est nullement une variété à bord !) $\hand{\partial}T$, qui admet
dans $T$ un voisinage isomorphe au produit de $T$ par un segment,
etc. Moyennant des hypothèses d'équisingularité convenables, on
s'attend à ce que $T$ soit muni, de façon essentiellement unique,
d'une structure de fibré localement trivial sur $Y$, admettant $\hand{\partial}T$
comme sous-fibré. C'est là un des points les moins clairs dans
l'intuition provisoire que j'ai de la situation, alors que la
classe d'homotopie de l'application structurale prévue $T \longrightarrow Y$
a un sens évident, indépendamment de toute hypothèse d'équisingula-%
rité, comme inverse homotopique de l'application d'inclusion $Y \longrightarrow T$,
qui doit être un homotopisme. Une façon d'obtenir a posteriori
une telle structure serait via l'hypothétique équivalence de caté-%
gories isotopiques envisagée au début, en tenant compte du fait que

\margin{111}

\clearpage
% Esquisse d'un programme (cote 119, copy 2): archive cote 119, site page 112 (PDF page 113), typescript page 34.
% typescript. Draft compared: esquisse/pages/esq_119_036.tex (copy 1, site 36).
% AI-transcribed by Claude Opus 5.5 (claude-opus-5-5, adaptive thinking, effort high), Message Batches API, batch esq119_full (msgbatch_015xqDHkTumDysRF1qngZhRJ), 2026-11-13; not yet checked by eye.
\scanpage{119/112}
\mspage{34}

le foncteur  $(U,Y,f) \longmapsto (X,Y)$   est défini de façon évidente,
indépendamment de toute théorie de voisinages tubulaire.

On dira sans doute, non sans quelque raison, que tout
celà n'est peut-être que rêves, qui s'évanouiront en fumée dès
qu'on s'essayera à un travail circonstancié, voire même dès avant
en face de certains faits connus ou bien évidents qui m'auraient
échappé. Certes, seul un travail sur pièces permettra de décan-%
ter le juste du faux et de connaître la substance véritable. La
seule chose dans tout cela qui ne fait pour moi l'objet d'aucun
doute, c'est la nécessité d'un tel travail de fondements, en
d'autres termes, la nature articielle des fondements actuels de
la topologie, et des difficultés que ceux-ci soulèvent à chaque
pas. Il est bien possible par contre que la formulation que je
donne à une théorie de dévissage des structures stratifiées, com-%
me un théorème d'équivalence de catégories isotopiques (voire même
$\infty$-isotopiques) convenables, soit trop optimiste. Je devrais ajou-%
ter pourtant que je n'ai guère de doutes non plus que la théorie
de ces dévissages que j'ai développé il y a deux ans, alors qu'elle
reste partiellement heuristique, exprime bel et bien une réalité
tout ce qu'il y a de palpable. Dans une partie de mon travail,
faute de pouvoir disposer d'un contexte ``modéré'' tout fait, et
pour avoir néanmoins des énoncés précis et démontrables, j'ai été
amené à postuler sur la structure stratifiée de départ des struc-%
tures supplémentaires tout ce qu'il y a de plausible, dans la na-%
ture de la donnée de rétractions locales notamment, qui dès lors
permettent bel et bien la construction d'un système canonique
d'espaces, paramétré par l'ensemble ordonné des ``drapeaux'' Drap(I)
de l'ensemble ordonné  $I$  indexant les strates, ces espaces jouant
le rôle des espaces  $(U,Y)$  de tantôt, reliés entre eux par des
applications de plongements et de fibrations propres, qui permet-%
tent de reconstituer de façon toute aussi canonique la structure
stratifiée de départ, y compris ces ``structures supplémentaires''(${}^{7}$).
Le seul ennui, c'est que ces dernières semblent un élément de struc-%
ture superfétatoire, qui n'est nullement une donnée dans les situations

\margin{112}

\clearpage
% Esquisse d'un programme (cote 119, copy 2): archive cote 119, site page 113 (PDF page 114), typescript page 35.
% typescript. Draft compared: esquisse/pages/esq_119_037.tex (copy 1, site 37).
% AI-transcribed by Claude Opus 5.5 (claude-opus-5-5, adaptive thinking, effort high), Message Batches API, batch esq119_full (msgbatch_015xqDHkTumDysRF1qngZhRJ), 2026-11-13; not yet checked by eye.
\scanpage{119/113}
\mspage{35}

géométriques courantes, par exemple pour l'espace modulaire com-%
pact $\hat{M}_{g,\nu}$ avec sa ``stratification à l'infini'' canonique, donnée
par le diviseur à croisements normaux de Mumford-Deligne. Une
autre difficulté, moins sérieuse sans doute, c'est que le soi-%
disant ``espace'' modulaire est en fait une \uline{multiplicité} - techni-%
quement, cela s'exprime surtout par la nécessité de remplacer
l'ensemble d'indices $I$ pour les strates par une \uline{catégorie} (es-%
sentiellement finie) d'indices, en l'occurence celle des
``graphes MD'', qui ``paramètrent'' les ``structures combinatoires''
possibles d'une courbe stable de type $(g,\nu)$. Ceci dit, je puis
affirmer que la théorie de dévissage général, spécialement déve-%
loppé sous la pression du besoin de \uline{cette} cause, s'est révélé
en effet un guide précieux, conduisant à une compréhension pro-%
gressive, d'une cohérence sans \unsure{failles,de|failles, de} certains aspects essen-%
tiels de la tour de Teichmüller (c'est à dire, essentiellement
de la``structure à l'infini'' des groupes de Teichmüller ordinaires).
C'est cette approche qui m'a conduit finalement, dans les mois
suivants, vers le principe d'une construction purement combina-%
toire de la tour des groupoïdes de Teichmüller, dans l'esprit es-%
quissé plus haut (cf. par. 2).

Un autre test de cohérence satisfaisant provient du
point de vue ``topossique''. En effet, mon intérêt pour les multi-%
plicités modulaires provenant avant tout de leur sens algébrico-%
géométrique et arithmétique, c'est aux multiplicités modulaires
\uline{algébriques}, sur le corps de base absolu $\mathbb{Q}$, que je me suis inté-%
ressé prioritairement, et à un ``dévissage'' à l'infini de leurs
groupes fondamentaux géométriques (i.e. des groupes de Teichmüller
\uline{profinis}) qui soit compatible avec les opérations naturelles de
$\mathrm{I}\!\Gamma = \mathrm{Gal}(\bar{\mathbb{Q}}/\mathbb{Q})$. Cela semblait exclure d'emblée la possibilité de me
référer à une hypothétique théorie de dévissage de structures
stratifiées dans un contexte de ``topologie modérée'' (ou même de
topologie ordinaire, cahin-caha), si ce n'est comme fil conducteur
entièrement heuristique. Dès lors se posait la question de tra-%
duire, dans le contexte des topos (en l'occurence les topos étales)

\margin{113}

\clearpage
% Esquisse d'un programme (cote 119, copy 2): archive cote 119, site page 114 (PDF page 115), typescript page 36.
% typescript. Draft compared: esquisse/pages/esq_119_038.tex (copy 1, site 38).
% AI-transcribed by Claude Opus 5.5 (claude-opus-5-5, adaptive thinking, effort high), Message Batches API, batch esq119_full (msgbatch_015xqDHkTumDysRF1qngZhRJ), 2026-11-13; not yet checked by eye.
\scanpage{119/114}
\mspage{36}

intervenant dans la situation, la théorie de dévissage à laquelle
j'étais parvenue dans un contexte tout différent - avec la tâche
supplémentaire, par la suite, de dégager un théorème de comparai-%
son général, sur le modèle des théorèmes bien connus, pour com-%
parer les invariants obtenus (notamment les types d'homotopie de
voisinages tubulaires divers) dans le cadre transcendant, et dans
le cadre schématique. J'ai pu me convaincre qu'un tel formalisme
de dévissage avait bel et bien un sens dans le contexte (dit
``abstrait'' !) des topos généraux, ou tout au moins des topos noe-%
thériens (comme ceux qui s'introduisent ici), via une notion conve-%
nable de \uline{voisinage tubulaire canonique d'un sous-topos} dans un
topos ambiant. Une fois cette notion acquise, avec certaines pro-%
priétés formelles simples, la description du ``dévissage'' d'un
topos stratifié est considérablement plus simple même dans ce
cadre, que dans le cadre topologique (modéré). Il est vrai que
là aussi il y a un travail de fondements à faire, notamment pour
la notion même de voisinage tubulaire d'un sous-topos - et il est
étonnant d'ailleurs que ce travail (pour autant que je sache)
n'ait toujours pas été fait, c'est-à-dire que personne (depuis plus
de vingt ans qu'il existe un contexte de topologies étale) ne
semble en avoir eu besoin ; un signe sûrement que la compréhension
de la structure topologique des schémas n'a pas tellement pro-%
gressée\deleted{s} depuis le travail d'Artin-Mazur ...

Une fois accompli le double travail de dégrossissage
(plus ou moins heuristique) autour de la notion de dévissage d'un
espace ou d'un topos stratifié, qui a été une étape cruciale dans
ma compréhension des multiplicités modulaires, il est d'ailleurs
apparu que pour les besoins de ces dernières, on peut sans doute
court-circuiter au moins une bonne partie de cette théorie par des
arguments géométriques directs. Il n'en reste pas moins que pour
moi, le formalisme de dévissage auquel je suis parvenu a fait ses
preuves d'utilité et de cohérence, indépendamment de toute ques-%
tion sur les fondements les plus adéquats qui permettent de lui
donner tout son sens.

\margin{114}

\clearpage
% Esquisse d'un programme (cote 119, copy 2): archive cote 119, site page 115 (PDF page 116), typescript page 37.
% typescript. Draft compared: esquisse/pages/esq_119_039.tex (copy 1, site 39).
% AI-transcribed by Claude Opus 5.5 (claude-opus-5-5, adaptive thinking, effort high), Message Batches API, batch esq119_full (msgbatch_015xqDHkTumDysRF1qngZhRJ), 2026-11-13; not yet checked by eye.
\scanpage{119/115}
\mspage{37}

6. Un des théorèmes de fondements de topologie (modérée)
les plus intéressants qu'il faudrait développer, serait un théo-%
rème de ``dévissage'' (encore !) d'une application modérée propre
d'espace modérés,

\[ f : X \longrightarrow Y , \]

via une filtration décroissante de $Y$ par des sous-espaces mo-%
dérés fermés $Y^{i}$, tels que au-dessus des ``strates ouvertes''
$Y^{i} \setminus Y^{i-1}$ de cette filtration, $f$ induise une fibration locale-%
ment triviale (du point de vue modéré, il va sans dire). Un tel
énoncé devrait encore se généraliser et se préciser de diverses
façons, notamment en demandant l'existence d'un dévissage ana-%
logue \uline{simultané}, pour $X$ et une famille finie donnée de sous-%
espaces (modérés) fermés de $X$. Egalement la notion même de fi-%
bration localement triviale au sens modéré peut se renforcer
considérablement, en tenant compte du fait que les strates ou-%
vertes $U_{i}$ sont \uline{mieux} que des espaces à structure modérée pu-%
rement locale, du fait qu'ils sont obtenus comme différence de
deux espaces modérés, compacts si $Y$ était compact. Entre la
notion d'espace modéré compact (qui se réalise comme un des
``modèles'' de départ dans un $\mathbb{R}^{n}$) et celle d'espace ``localement
modéré'' (localement compact) qui s'en déduit de façon assez évi-%
dente, il y a une notion un peu plus délicate d'\uline{espace} \uline{``globa-}%
\uline{lement modéré}'' $X$, obtenu comme différence $X^{\wedge}\setminus Y$ de deux espaces
modérés compacts, étant entendu qu'on ne distingue pas entre
l'espace défini par une paire $(X^{\wedge},Y)$, et celui défini par une
paire $(X'^{\wedge},Y')$ qui s'en déduit par une application modérée
(nécessairement propre)

\[ g : X'^{\wedge} \longrightarrow X^{\wedge} \]

induisant une bijection $g^{-1}(X) \xrightarrow{\ \sim\ } X$, en prenant $Y' = g^{-1}(Y)$.
L'exemple naturel le plus intéressant peut-être est celui où on
part d'un schéma séparé de type fini sur $\mathbb{C}$ ou sur $\mathbb{R}$, en prenant

\margin{115}

\clearpage
% Esquisse d'un programme (cote 119, copy 2): archive cote 119, site page 116 (PDF page 117), typescript page 38.
% typescript. Draft compared: esquisse/pages/esq_119_040.tex (copy 1, site 40).
% AI-transcribed by Claude Opus 5.5 (claude-opus-5-5, adaptive thinking, effort high), Message Batches API, batch esq119_full (msgbatch_015xqDHkTumDysRF1qngZhRJ), 2026-11-13; not yet checked by eye.
\scanpage{119/116}
\mspage{38}

pour  $X$  l'ensemble de ses points complexes ou réels, qui hérite
d'une structure modérée globale à l'aide des compactifications
schématique (qui existent d'après Nagata) du schéma de départ.
Cette notion d'espace globalement modéré est associée à une no-%
tion d'\uline{application globalement modérée}, qui permet à son tour de
renforcer en conséquence la notion de fibration localement tri-%
viale, dans l'énoncé d'un théorème de dévissage pour une applica-%
tion  $f : X \longrightarrow Y$  (pas nécessairement propre maintenant) dans
le contexte des espaces globalement modérés.

J'ai été informé l'été dernier par Zoghman Mebkhout
qu'un théorème de dévissage dans cet esprit avait été obtenu ré-%
cemment dans le contexte des espaces analytiques réels et/ou com-%
plexes, avec des  $Y^{i}$  qui, cette fois, sont des sous-espaces ana-%
lytiques de  $Y$. Ce résultat rend plausible qu'on dispose dès à
présent de moyens techniques suffisamment puissants pour démontrer
également un théorème de dévissage dans le contexte modéré, plus
général en apparence, mais probablement moins ardu.

C'est le contexte d'une topologie modérée également
qui devrait permettre, il me semble, de formuler avec précision
un principe général très sûr que j'utilise depuis longtemps dans
un grand nombre de situations géométriques, que j'appelle le
``\uline{principe des choix anodins}'' - aussi utile que vague d'apparence !
Il dit, lorsque pour les besoins d'une construction quelconque
d'un objet géométrique en termes d'autres, on est amené à faire
un certain nombre de choix arbitraires en cours de route, de façon
donc que l'objet obtenu dépend en apparence de ces choix, et est
donc entâché d'un défaut de canonicité, que ce défaut est sérieux
en effet (et pour être levé demande une analyse plus soigneuse de
la situation, des notions utilisées, des données introduites etc )
chaque fois que l'un au moins de ces choix s'effectue dans un
``espace'' qui n'est pas ``contractile'' i.e. dont le  $\pi_{0}$  \uline{ou} un des
invariants supérieurs  $\pi_{i}$  est non trivial ; que ce défaut est
par contre apparent seulement, que la construction est ``essentiel-%
lement canonique'' et n'entraînera pas vraiment d'ennuis, chaque

\margin{116}

\clearpage
% Esquisse d'un programme (cote 119, copy 2): archive cote 119, site page 117 (PDF page 118), typescript page 39.
% typescript. Draft compared: esquisse/pages/esq_119_041.tex (copy 1, site 41).
% AI-transcribed by Claude Opus 5.5 (claude-opus-5-5, adaptive thinking, effort high), Message Batches API, batch esq119_full (msgbatch_015xqDHkTumDysRF1qngZhRJ), 2026-11-13; not yet checked by eye.
\scanpage{119/117}
\mspage{39}

fois que les choix faits sont tous ``anodins'' i.e. s'effectuent
dans des espaces \uline{contractiles}. Quand on essaye dans les cas d'es-%
pèce de cerner de plus près ce principe, il semble qu'on tombe à
chaque fois sur la notion de ``catégories $\infty$- isotopiques'' exprimant
une situation donnée, plus fines que les catégories isotopiques
(= 0-isotopiques) plus naïves, obtenues en ne retenant que les $\pi_0$
des espaces d'isomorphismes qui s'introduisent dans la situation,
alors que le point de vue $\infty$ -isotopique retient tout leur type
d'homotopie. Par exemple, le point de vue isotopique naïf pour
les surfaces compactes à bord orientées de type $(g,\nu)$ est ``bon''
(sans boomerang caché !) exactement dans les cas que j'appelle
``anabéliens'' (et que Thurston appelle ``hyperboliques'') i.e. dis-%
tincts de $(0,0)$, $(0,1)$, $(0,2)$, $(1,0)$ - qui sont aussi les cas
justement où le groupe des automorphismes de la surface a une
composante neutre \uline{contractile}. Dans les autres cas, sauf le cas
$(0,0)$ de la sphère sans trou, il suffit de travailler avec les
catégories 1-isotopiques pour exprimer de façon satisfaisante par
voie algébrique les faits géométrico-topologiques essentiels, vu
que ladite composante connexe est alors un $K(\pi,1)$. Travailler
dans une catégorie 1-isotopique revient d'ailleurs à travailler
dans une bicatégorie, i.e. avec des $\underline{\mathrm{Hom}}(X,Y)$ qui sont (non plus
des ensembles discrets comme dans le point de vue 0-isotopique,
mais) des groupoïdes (dont les $\pi_0$ ne sont autres que les Hom
0-isotopiques). C'est la description en termes purement algébri-%
ques de cette bicatégorie qui est faite dans la dernière partie
de la thèse de Yves Ladegaillerie (cf. par. 3).

Si je me suis étendu ici plus longuement sur le thème
des fondements de la topologie modérée, qui n'est nullement un de
ceux auxquels je compte me consacrer prioritairement dans les an-%
nées qui viennent, c'est sans doute justement que je sens qu'il
y a là d'autant plus une cause qui a besoin d'être plaidée, ou
plutôt : un travail d'une grande actualité qui a besoin de bras !
Comme naguère pour de nouveaux fondements de la géométrie algé-%
brique, ce ne sont pas des plaidoyers qui surmontent l'inertie

\margin{117}

\clearpage
% Esquisse d'un programme (cote 119, copy 2): archive cote 119, site page 118 (PDF page 119), typescript page 40.
% typescript. Draft compared: esquisse/pages/esq_119_042.tex (copy 1, site 42).
% AI-transcribed by Claude Opus 5.5 (claude-opus-5-5, adaptive thinking, effort high), Message Batches API, batch esq119_full (msgbatch_015xqDHkTumDysRF1qngZhRJ), 2026-11-13; not yet checked by eye.
\scanpage{119/118}
\mspage{40}

des habitudes acquises, mais un travail tenace, méticuleux, sans
doute de longue haleine, et porteur au jour le jour de moissons
éloquentes.

\margin{\unsure{espacer}}Je voudrais encore dire quelques mots sur une réflexion
plus ancienne (fin des années 60 ?), très proche de celle dont il
vient d'être question, inspirée par les idées de Nash, qui m'avaient
beaucoup frappé. Au lieu ici de définir axiomatiquement une notion
de ``théorie modérée'' via la donnée d'une notion de ``partie modérée
de $\mathbb{R}^n$ '' satisfaisant à certaines conditions (de stabilité surtout),
c'est à une axiomatisation de la notion de ``variété lisse'' et du
formalisme différentiable sur de telles variétés que j'en avais,
via la donnée, pour chaque entier naturel $n$, d'un sous-anneau
$\mathbb{A}_n$ de l'anneau des germes de fonctions réelles à l'origine dans $\mathbb{R}^n$.
Ce sont les fonctions qui seront admises pour exprimer les ``change-%
ments de carte'' pour la notion de $\mathbb{A}$-variété correspondante, et il
s'est agit de dégager tout d'abord un système d'axiomes sur le
système $\mathbb{A} = (\mathbb{A}_n)_{n\in\underline{N}}$ qui assure à cette notion de variété une sou-%
plesse comparable à celle de variété $C^\infty$, ou analytique réelle
(ou de Nash). Suivant le type de constructions familières qu'on
tient à pouvoir effectuer dans le contexte des $\mathbb{A}$-variétés, le
système d'axiomes pertinent est plus ou moins réduit, ou riche.
Très peu suffit s'il s'agit seulement de développer le formalisme
différentiel, avec la construction de fibrés de jets, les complexes
de De Rham etc. Si on veut un énoncé du type ``quasi-fini implique
fini'' (pour une application au voisinage d'un point), qui est ap-%
paru comme un énoncé-clef dans la théorie locale des espaces ana-%
lytiques, il faut un axiome de stabilité de nature plus délicate,
dans le ``Vorbereitungssatz'' de Weierstrass (*). Dans d'autres

\rule{4cm}{0.4pt}

(*) Il peut paraître plus simple de dire que les anneaux (locaux)
$\mathbb{A}_n$ sont \uline{henséliens}, ce qui est équivalent. Mais il n'est nullement
clair a priori sous cette dernière forme que la condition en ques-%
tion est dans la nature d'une condition de stabilité, circonstance
importante comme il apparaîtra dans les réflexions qui suivent.

\margin{118}

\clearpage
% Esquisse d'un programme (cote 119, copy 2): archive cote 119, site page 119 (PDF page 120), typescript page 41.
% typescript. Draft compared: esquisse/pages/esq_119_043.tex (copy 1, site 43).
% AI-transcribed by Claude Opus 5.5 (claude-opus-5-5, adaptive thinking, effort high), Message Batches API, batch esq119_full (msgbatch_015xqDHkTumDysRF1qngZhRJ), 2026-11-13; not yet checked by eye.
\scanpage{119/119}
\mspage{41}

questions, un axiome de stabilité par prolongement analytique
(dans $\mathbb{C}^{n}$) apparaît nécessaire. L'axiome le plus draconien que
j'aie été amené à introduire, lui aussi un axiome de stabilité,
concerne l'intégration des systèmes de Pfaff, assurant que cer-%
tains groupes de Lie, voire tous, sont des $\mathcal{A}$-variétés. Dans tout
ceci, \uline{j'ai pris soin de ne pas supposer que les} $\mathcal{A}_{n}$ \uline{soient des}
\uline{$\mathbb{R}$-algèbres}, donc une fonction constante sur une $\mathcal{A}$-variété n'est
``admissible'' que si sa valeur appartient à un certain sous-corps
$K$ de $\mathbb{R}$ (c'est , si on veut, $\mathcal{A}_{0}$). Ce sous-corps peut fort bien
être $\mathbb{Q}$, ou sa fermeture algébrique $\bar{\mathbb{Q}}_{r}$ dans $\mathbb{R}$, ou toute autre
sous-extension de $\mathbb{R}/\mathbb{Q}$ , de préférence même de degré de trans-%
cendance fini, ou du moins dénombrable, sur $\mathbb{Q}$. Cela permet par
exemple, comme tantôt pour les espaces modérés, de faire corres-%
pondre à tout point $x$ d'une variété (de type $\mathcal{A}$) un corps rési-%
dual $k(x)$, qui est une sous-extension de $\mathbb{R}/K$. Un fait qui me
semble important ici, c'est que même sous sa forme la plus forte,
le système d'axiomes n'implique \uline{pas} qu'on doive avoir $K = \mathbb{R}$. Plus
précisément, du fait que \uline{tous} les axiomes sont des axiomes de
stabilité, il résulte que pour un ensemble $S$ donné de germes
de fonctions analytiques réelles à l'origine (dans divers expaces
$\mathbb{R}^{n}$), il existe une plus petite théorie $\mathcal{A}$ pour laquelle ces germes
sont admissibles, et que celle-ci est ``dénombrable'' i.e. les $\mathcal{A}_{n}$
sont dénombrables, dès que $S$ l'est. A fortiori, $K$ est alors dé-%
\margin{x} nombrable, i.e. de degré de transcendance dénombrable sur $\mathbb{Q}$.

L'idée ici est d'introduire, par le biais de cette
axiomatique, une notion de fonction (analytique réelle) ``\uline{élémentaire}'',
ou plutôt, toute une hiérarchie de telles notions. Pour une fonc-%
tion de 0 variables i.e. une constante, cette notion donne celle
de ``constante élémentaire'', incluant notamment (dans le cas de
l'axiomatique la plus forte) des constantes telles que $\pi$, $e$ et
une multitude d'autres, en prenant des valeurs de fonctions admis-%
sibles (telles l'exponentielle, le logarithme etc) pour des systèmes
de valeurs ``admissibles'' de l'argument. On sent que la relation
entre le système $\mathcal{A} = (\mathcal{A}_{n})_{n\in\underline{N}}$ et le corps de rationalité $K$

\margin{119}

\clearpage
% Esquisse d'un programme (cote 119, copy 2): archive cote 119, site page 120 (PDF page 121), typescript page 42.
% typescript. Draft compared: esquisse/pages/esq_119_044.tex (copy 1, site 44).
% AI-transcribed by Claude Opus 5.5 (claude-opus-5-5, adaptive thinking, effort high), Message Batches API, batch esq119_full (msgbatch_015xqDHkTumDysRF1qngZhRJ), 2026-11-13; not yet checked by eye.
\scanpage{119/120}
\mspage{42}

correspondant doit être très étroite, du moins pour des $\mathcal{A}$ qui
peuvent être engendrés par un ``système de générateurs'' S fini -
mais il est à craindre que la moindre question intéressante qu'on
pourrait se poser sur cette situation soit actuellement hors de
portée ($^1$).

Ces réflexions anciennes ont repris quelque actualité
pour moi avec ma réflexion ultérieure sur les théories modérées.
Il me semble en effet qu'il est possible d'associer de façon na-%
turelle à une ``théorie différentiable'' $\mathcal{A}$ une théorie modérée $\mathcal{T}$
(ayant sans doute même corps de constantes), de telle façon que
tout $\mathcal{A}$-variété soit automatiquement munie d'une structure $\mathcal{T}$-modérée,
et inversement que pour tout espace $\mathcal{T}$-modéré compact $X$, on puisse
trouver une partie fermée modérée rare $Y$ dans $X$, telle que $X \smallsetminus Y$
provienne d'une $\mathcal{A}$-variété, et que de plus cette structure de $\mathcal{A}$-%
variété soit unique tout au moins dans le sens suivant : deux tel-%
les structures coïncident dans le complémentaire d'une partie modé-%
rée rare $Y' \supset Y$ de $X$. La théorie de dévissage des structures
modérées stratifiées (dont il a été question au par. précédent),
dans le cas de strates lisses, devrait d'ailleurs soulever des
questions beaucoup plus précises encore de comparaison des struc-%
tures modérées avec des structures de type différentiable (ou plu-%
tôt, $\mathbb{R}$-analytique). Je soupçonne que le type d'axiomatisation pro-%
posé ici pour la notion de ``théorie différentiable'' fournirait un
cadre naturel pour formuler de telles questions avec toute la pré-%
cision et toute la généralité souhaitables.


7. Depuis le mois de mars de l'an dernier, donc depuis près
d'un an, la plus grande partie de mon énergie a été consacrée à un
travail de réflexion sur les \uline{fondements de l'algèbre} (\uline{co})\uline{homologique}
\uline{non commutative}, ou ce qui revient au même, finalement, de l'\uline{algèbre}
\uline{homotopique}. Ces réflexions se sont concrétisées par un volumineux
paquet de notes dactylographiées, destinées à former le premier
volume (actuellement en cours d'achèvement) d'un ouvrage en deux

\margin{120}

\clearpage
% Esquisse d'un programme (cote 119, copy 2): archive cote 119, site page 121 (PDF page 122), typescript page 43.
% typescript. Draft compared: esquisse/pages/esq_119_045.tex (copy 1, site 45).
% AI-transcribed by Claude Opus 5.5 (claude-opus-5-5, adaptive thinking, effort high), Message Batches API, batch esq119_full (msgbatch_015xqDHkTumDysRF1qngZhRJ), 2026-11-13; not yet checked by eye.
\scanpage{119/121}
\mspage{43}

volumes à paraître chez Hermann, sous le titre commun ``\uline{A la Poursuite}
\uline{des Champs}''. Je prévois actuellement (après des élargissements
successifs du propos initial) que le manuscript de l'ensemble des
deux volumes, que j'espère achever en cours d'année pour ne plus
avoir à y revenir, fera dans les 1 500 pages dactylographiées.
Ces deux volumes d'ailleurs sont pour moi les premiers dans une
série plus vaste, sous le titre commun ``\uline{Réflexions Mathématiques}'',
où je compte développer tant soit peu certains des thèmes esquis-%
sés dans le présent rapport.

Vu qu'il s'agit d'un travail en cours de rédaction,
et même d'achèvement, dont le premier volume sans doute paraîtra
cette année et contiendra une introduction circonstanciée, il
est sans doute moins intéressant que je m'étende ici sur ce thème
de réflexion, et me contenterai donc d'en parler très brièvement.
Ce travail me semble quelque peu marginal par rapport aux thèmes
que je viens d'esquisser, et ne représente pas (il me semble) un
véritable renouvellement d'optique ou d'approche par rapport à
mes intérêts et ma vision mathématiques d'avant 1970. Si je m'y
suis résolu soudain, c'est presque en dése\over{r}{s}poir de cause, alors
que près de vingt ans se sont écoulés depuis que se sont posées
en termes bien clairs un certain nombre de questions visiblement
fondamentales, et mûres pour être menées à leur terme, sans que
personne ne les voie, ou prenne la peine de les sonder. Aujour-%
d'hui encore, les structures de base qui interviennent dans le
point de vue homotopique en topologie, y compris même en algèbre
homologique commutative, ne sont pas comprises, et à ma connais-%
sance, après les travaux de Verdier, de Giraud et d'Illusie sur
ce thème (qui constituent autant de ``coups d'envoi'' attendant tou-%
jours une suite ...) il n'y a pas eu d'effort dans ce sens. Je
devrais faire exception sans doute pour le travail d'axiomatisa-%
tion fait par Quillen sur la notion de catégorie de modèles, à
la fin des années 60, et repris sous des variantes diverses par
divers auteurs. Ce travail à l'époque, et maintenant encore, m'a
beaucoup séduit et appris, tout en allant dans une direction assez

\margin{121}

\clearpage
% Esquisse d'un programme (cote 119, copy 2): archive cote 119, site page 122 (PDF page 123), typescript page 44.
% typescript. Draft compared: esquisse/pages/esq_119_046.tex (copy 1, site 46).
% AI-transcribed by Claude Opus 5.5 (claude-opus-5-5, adaptive thinking, effort high), Message Batches API, batch esq119_full (msgbatch_015xqDHkTumDysRF1qngZhRJ), 2026-11-13; not yet checked by eye.
\scanpage{119/122}
\mspage{44}

différente de celle qui me tenait et tient à coeur. Il introduit
certes des catégories dérivées dans divers contextes non commu-%
tatifs, mais sans entrer dans la question des structures internes
essentielles d'une telle catégorie, laissée ouverte également
dans le cas commutatif par Verdier, et après lui par Illusie. De
même, la question de mettre le doigt sur les ``coefficients''
naturels pour un formalisme cohomologique non commutatif, au delà
des champs (qu'on devrait appeler 1-champs) étudiés dans le livre
de Giraud, restait ouverte - ou plutôt, les intuitions riches et
précises qui y répondent, puisées dans des exemples nombreux pro-%
venant de la géométrie algébrique notamment, attendent toujours
un langage précis et souple pour leur donner forme.

Je reviens sur certains aspects de ces questions de
fondements en 1975, à l'occasion (je crois me souvenir) d'une
correspondance avec Larry Breen (\deleted{deux}\inserted{trois} lettres de cette corres-%
pondance seront reproduites en appendice au Chap. I du volume 1,
``Histoire de Modèles'', de la Poursuite des Champs). A ce moment
apparaît l'intuition que les $\infty$-groupoïdes doivent constituer des
modèles, particulièrement adéquats, pour les types d'homotopie,
les n-groupoïdes correspondant aux types d'homotopie \uline{tronqués}
(avec $\pi_i = 0$ pour $i > n$). Cette même intuition, par des voies
très différentes, a été retrouvée par Ronnie Brown à Bangor et
certains de ses élèves, mais en utilisant une notion de $\infty$-groupoïde
assez restrictive (qui, parmi les types d'homotopie 1-connexes,
ne modélise que les produits d'espacesd'Eilenberg-Maclane). C'est
stimulé par une correspondance à bâtons rompus avec Ronnie Brown,
que j'ai finalement repris une réflexion, commençant par un essai
de définition d'une notion de $\infty$-groupoïde plus large (rebaptisé
par la suite ``champ en $\infty$-groupoïdes'' ou simplement ``champ'', sous-%
entendu : sur le topos ponctuel), et qui de fil en aiguille m'a
amené à la Poursuite des Champs. Le volume ``Histoire de Modèles''
y constitue d'ailleurs une digression entièrement imprévue par
rapport au propos initial (les fameux champs étant provisoirement
oubliés, et n'étant prévus réapparaître que vers les pages 1 000
environ ...).

\margin{122}

\clearpage
% Esquisse d'un programme (cote 119, copy 2): archive cote 119, site page 123 (PDF page 124), typescript page 45.
% typescript. Draft compared: esquisse/pages/esq_119_047.tex (copy 1, site 47).
% AI-transcribed by Claude Opus 5.5 (claude-opus-5-5, adaptive thinking, effort high), Message Batches API, batch esq119_full (msgbatch_015xqDHkTumDysRF1qngZhRJ), 2026-11-13; not yet checked by eye.
\scanpage{119/123}
\mspage{45}

Ce travail n'est pas entièrement isolé par rapport à
mes intérêts plus récents. Par exemple, ma réflexion sur les
multiplicités modulaires $\widehat{M}_{g,\nu}$ et leur structure stratifiée
a relancé une réflexion sur un théorème de van Kampen de di-%
mension $>$ 1 (un des thèmes de prédilection également du groupe
de Bangor), et a peut-être contribué à préparer le terrain pour
le travail de plus grande envergure l'année d'après. Celui-ci
rejoint également par moments une réflexion datant de la même
année 1975 (ou l'année d'après) sur un ``complexe de De Rham à
puissances divisées'', qui a fait l'objet de ma dernière confé-%
rence publique, à l'IHES en 1976, et dont le manuscrit, confié
je ne me rappelle plus à qui après l'exposé, est d'ailleurs perdu.
C'est au moment de cette réflexion que germe aussi l'intuition
d'une ``schématisation'' des types d'homotopie, que sept ans après
j'essaye de préciser dans un chapître (particulièrement hypothé-%
tique) de l'Histoire de Modèles.

Le travail de réflexion entrepris dans la Poursuite
des Champs est un peu comme une dette dont je m'acquitterais,
vis-à-vis d'un passé scientifique où, pendant une quinzaine
d'années (entre 1955 et 1970), le développement d'outils cohomo-%
logiques a été le Leitmotiv constant, dans mon travail de fonde-%
ments de la géométrie algébrique. Si la reprise actuelle de ce
thème-là a pris des dimensions inattendues, ce n'est pas cepen-%
dant par pi\over{ti}{\unsure{ét}}é pour un passé, mais à cause des nombreux imprévus
faisant irruption sans cesse, en bousculant sans ménagement les
plans et propos prévus - un peu comme dans un conte des mille et
une nuits, où l'attention se trouve maintenue en haleine à tra-%
vers vingt autres contes avant de connaître le fin mot du premier.


8. J'ai très peu parlé encore des réflexions plus terre
à terre de géométrie topologique bidimensionelle, associées no-%
tamment à mes activités d'enseignant et celles dites de ``direc-%
tion de recherches''. A plusieurs reprises, j'ai vu s'ouvrir devant

\margin{123}

\clearpage
% Esquisse d'un programme (cote 119, copy 2): archive cote 119, site page 124 (PDF page 125), typescript page 46.
% typescript. Draft compared: esquisse/pages/esq_119_048.tex (copy 1, site 48).
% AI-transcribed by Claude Opus 5.5 (claude-opus-5-5, adaptive thinking, effort high), Message Batches API, batch esq119_full (msgbatch_015xqDHkTumDysRF1qngZhRJ), 2026-11-13; not yet checked by eye.
\scanpage{119/124}
\mspage{46}

moi de vastes et riches champs mûrs pour la moisson, sans que
jamais je réussisse à communiquer cette vision, et l'étincelle
qui l'accompagne, à un (ou une) de mes élèves, et à la faire
déboucher sur une exploration commune, de plus ou moins longue
haleine. A chaque fois jusqu'à aujourd'hui même, après une pros-%
pection de quelques jours ou quelques semaines, où je découvrais
en éclaireur des richesses insoupçonnées au départ, le voyage
tournait court, quand il devenait clair que je serai seul à le
poursuivre. Des intérêts plus forts prenaient le pas alors sur
un voyage qui, dès lors, apparaissait comme une digression, voire
une dispersion, plutôt qu'une aventure poursuivie en commun.

Un de ces thèmes a été celui des polygones plans, centré
autour des variétés modulaires qu'on peut leur associer. Une des
surprises ici a été l'irruption de la géométrie algébrique dans
un contexte qui m'en avait semblé bien éloigné. Ce genre de sur-%
prise, lié à l'ubiquité de la géométrie \inserted{algébrique} dans la géométrie tout
court, s'est d'ailleurs répé\over{r}{t}é à plusieurs reprises.

Un autre thème a été celui des courbes (notamment des
cercles) immergés dans une surface, avec une attention particu-%
lière pour le cas ``stable'' où les points singuliers sont des points
doubles ordinaires (et aussi celui, plus général, où les diffé-%
rentes branches en un point se croisent mutuellement), avec sou-%
vent l'hypothèse supplémentaire que l'immersion soit ``cellulaire'',
i.e. donne naissance à une carte. Une variante de situations de
ce type est celle des immersions d'une surface à bord non vide,
et en tout premier lieu d'un disque (qui m'avait été signalé par
A. Campo il y a une dizaine d'années). Au delà de la question de
diverses formulations combinatoires de telles situations, qui ne
représente plus guère qu'un exercice de syntaxe, je me suis inté-%
ressé surtout à une vision dynamique des configurations possibles,
avec le passage de l'une à l'autre par déformations continues, qui
peuvent se décomposer en composées de deux types d'\uline{opérations élé-}%
\uline{mentaires} et leurs inverses, savoir le ``\uline{balayage}'' d'une branche
de courbe par dessus un point double, et l'\uline{effacement} ou \uline{la création}

\margin{124}

\clearpage
% Esquisse d'un programme (cote 119, copy 2): archive cote 119, site page 125 (PDF page 126), typescript page 47.
% typescript. Draft compared: esquisse/pages/esq_119_049.tex (copy 1, site 49).
% AI-transcribed by Claude Opus 5.5 (claude-opus-5-5, adaptive thinking, effort high), Message Batches API, batch esq119_full (msgbatch_015xqDHkTumDysRF1qngZhRJ), 2026-11-13; not yet checked by eye.
\scanpage{119/125}
\mspage{47}

\uline{d'un bigône}. (La première de ces opérations joue également un
rôle-clef dans une théorie ``dynamique'' des systèmes de pseudo-%
droites dans un plan projectif réel.) Une des premières questions
qui se posent ici est celle de déterminer les différentes \uline{classes
d'immersions} d'un cercle ou d'un disque (disons) modulo ces opé-%
rations élémentaires ; une autre, celle de voir quelles sont les
immersions du bord du disque qui proviennent d'une immersion du
disque, et dans quelle mesure les premières déterminent les se-%
condes. Ici encore, il m'a semblé que c'est une étude systématique
des variétés modulaires pertinentes (de dimension infinie en l'oc-%
curence, à moins d'arriver à en donner une version purement com-%
binatoire) qui devrait fournir le ``focus'' le plus efficace, nous
forçant en quelque sorte à nous poser les questions les plus per-%
tinentes. Malheureusement, la réflexion sur les questions même
les plus évidentes et \inserted{les} plus terre-à-terre est restée à l'état em-%
bryonnaire. Comme seul résultat tangible, je peux signaler une
théorie de ``dévissage'' canonique d'une immersion cellulaire sta-%
ble du cercle dans une surface, en immersions ``indécomposables'',
par ``téléscopage'' de telles immersions. Je n'ai pas réussi malheu-%
reusement à voir se transformer mes lumières sur la question en
un travail de stage de DEA, ni d'autres lumières (sur une descrip-%
tion théorique complète, en termes de groupes fondamentaux de
1-complexes topologiques, des immersions d'une surface à bord qui
prolongent une immersion donnée de son bord) en le démarrage d'une
thèse de doctorat d'état...

Un troisième thème, poursuivi simultanément depuis trois
ans à divers niveaux d'enseignement (depuis l'option pour étudiants
de première année, jusqu'à trois thèses de troisième cycle actuel-%
lement poursuivies sur ce thème) porte sur la classification to-%
pologique-combinatoire des systèmes de droites ou pseudo-droites.
Dans l'ensemble, la participation de mes élèves ici a été moins
décevante qu'ailleurs, et j'ai eu le plaisir parfois d'apprendre
par eux des choses intéressantes auxquelles je n'aurais pas songé.
La réflexion commune, par la force des choses, s'est limitée cepen-%

\margin{125}

\clearpage
% Esquisse d'un programme (cote 119, copy 2): archive cote 119, site page 126 (PDF page 127), typescript page 48.
% typescript. Draft compared: esquisse/pages/esq_119_050.tex (copy 1, site 50).
% AI-transcribed by Claude Opus 5.5 (claude-opus-5-5, adaptive thinking, effort high), Message Batches API, batch esq119_full (msgbatch_015xqDHkTumDysRF1qngZhRJ), 2026-11-13; not yet checked by eye.
\scanpage{119/126}
\mspage{48}

dant à un niveau très élémentaire. Dernièrement, j'ai finalement,
consacré un mois de réflexion intensive au développement d'une
construction purement combinatoire d'une sorte de ``surface modulaire''
associée à un système de $n$ pseudo-droites, qui classifie les
différentes ``positions relatives'' possibles (stables ou non) d'une
$(n+1)$ème pseudo-droite par rapport au système donné, ou encore :
les différentes ``affinisations'' possibles de ce système, par les
différents choix possibles d'une ``pseudo-droite à l'infini''. J'ai
l'impression d'avoir mis le doigt sur un objet remarquable, fai-%
sant apparaître un ordre imprévu dans des questions de classifi-%
cation qui jusqu'à présent apparaissaient assez chaotiques !
Mais ce n'est pas le lieu dans le présent rapport de m'étendre plus
à ce sujet.

Depuis 1977, dans toutes les questions (comme dans ces
deux derniers thèmes que je viens d'évoquer) où interviennent des
cartes bidimensionelles, la possibilité de les réaliser canoni-%
quement sur une surface conforme, donc sur une courbe algébrique
complexe dans le cas orienté compact, reste en filigrane constant
dans ma réflexion. Dans pratiquement tous les cas (en fait, tous
les cas sauf celui de certaines cartes sphériques avec ``peu d'auto-%
morphismes'') une telle réalisation conforme implique en fait une
\uline{métrique riemanienne canonique}, ou du moins, canonique à une cons-%
tante multiplicative près. Ces nouveaux éléments de structure
(sans même prendre en compte l'élément arithmétique, dont il a
été question au par. 3) sont de nature à transformer profondément
l'aspect initial des questions abordées, et les méthodes d'appro-%
che. Un début de familiarisation avec les belles idées de Thurston
sur la construction de l'espace de Teichmüller, en termes d'un jeu
très simple de chirurgie riemanienne hyperbolique, me confirme
dans ce pressentiment. Malheureusement, le niveau de culture très
modeste de presque tous les élèves qui ont travaillé avec moi pen-%
dans ces dix dernières années ne me permet pas d'aborder avec eux,
ne serait-ce que par allusion, de telles possibilités, alors que
l'assimilation d'un langage combinatoire minimum se heurte déjà,

\margin{126}

\clearpage
% Esquisse d'un programme (cote 119, copy 2): archive cote 119, site page 127 (PDF page 128), typescript page 49.
% typescript. Draft compared: esquisse/pages/esq_119_051.tex (copy 1, site 51).
% AI-transcribed by Claude Opus 5.5 (claude-opus-5-5, adaptive thinking, effort high), Message Batches API, batch esq119_full (msgbatch_015xqDHkTumDysRF1qngZhRJ), 2026-11-13; not yet checked by eye.
\scanpage{119/127}
\mspage{49}

bien souvent, à des obstacles psychiques considérables. C'est pour-%
quoi, à certains égards et de plus en plus ces dernières années,
mes activités d'enseignant ont souvent agi comme un poids, plu-%
tôt que comme un stimulant pour le déploiement d'une réflexion
géométrique tant soit peu avancée, ou seulement délicate.

9. L'occasion me semble propice ici de faire un bref bi-%
lan de mon activité enseignante depuis 1970, c'est-à-dire depuis
que celle-ci s'effectue dans un cadre universitaire. Ce contact
avec une réalité très différente a été pour moi riche en ensei-%
gnements, d'une portée d'un tout autre ordre d'ailleurs que sim-%
plement pédagogique ou scientifique. Ce n'est pas ici le lieu de
m'étendre sur ce sujet. J'ai dit aussi au début de ce rapport le
rôle qu'a joué ce changement de milieu professionnel dans le
renouvellement de mon approche des mathématiques, et celui de mes
centres d'intérêt en mathématique. Si par contre je fais le bilan
de mon activité enseignante au niveau de la recherche proprement
dite, j'aboutis à un constat d'échec clair et net. Depuis plus
de dix ans que cette activité se poursuit an par an au sein d'une
même institution universitaire, je n'ai pas su, à aucun moment,
y susciter un lieu où ``il se passe quelque chose'' - où quelque
chose ``passe'', parmi un groupe si réduit soit-il de personnes, re-%
liées par une aventure commune. A deux reprises, il est vrai, vers
les années 74 à 76, j'ai eu le plaisir et le privilège de susciter
chez un élève un travail d'envergure, poursuivi avec élan : chez
Yves Ladegaillerie le travail signalé précédemment (par. 3) sur
les questions d'isotopie en dimension 2, et chez Carlos Contou-Carrère
(dont la passion mathématique n'avait pas attendu la rencontre
avec moi pour éclore) un travail non publié sur les jacobiennes
locales et globales sur des schémas de bases généraux (dont une
partie a été annoncée dans une note aux CR). Ces deux cas mis à
part, mon rôle s'est borné, au cours de ces dix ans, à transmettre
tant bien que mal des rudiments du métier de mathématicien\deleted{\unsure{s}} à quel-%
ques vingt élèves au niveau de la recherche, ou tout au moins à

\margin{127}

\clearpage
% Esquisse d'un programme (cote 119, copy 2): archive cote 119, site page 128 (PDF page 129), typescript page 50.
% typescript. Draft compared: esquisse/pages/esq_119_052.tex (copy 1, site 52).
% AI-transcribed by Claude Opus 5.5 (claude-opus-5-5, adaptive thinking, effort high), Message Batches API, batch esq119_full (msgbatch_015xqDHkTumDysRF1qngZhRJ), 2026-11-13; not yet checked by eye.
\scanpage{119/128}
\mspage{50}

ceux parmi eux qui ont persévéré suffisamment avec moi, réputé
plus exigeant que d'autres, pour aboutir à un premier travail
noir sur blanc acceptable (certaines fois aussi à un travail mieux
qu'acceptable et plus qu'\over{e}{\unsure{u}}n seul travail, fait avec goût et jus-%
qu'au bout). Vu la conjoncture, même parmi les rares qui ont per-%
sévéré, plus rares encore seront ceux qui auront l'occasion
d'exercer ce métier, et par là, tout en gagnant leur pain, de
l'approfondir.

10. Depuis l'an dernier, je sens qu'au cours de mon acti-%
vité d'enseignant universitaire, j'ai appris tout ce que j'avais
à en apprendre et enseigné tout ce que je peux y enseigner, et
qu'elle a cessé d'être vraiment utile, à moi-même comme aux au-%
tres. M'obstiner sous ces conditions à la poursuivre encore me
paraîtrait un gaspillage, tant de ressources humaines que de deniers
publics. C'est pourquoi j'ai demandé mon détachement au CNRS (que
j'avais quitté en 1959 comme directeur de recherches frais émoulu,
pour entrer à l'IHES). Je sais d'ailleurs que la situation de
l'emploi est serrée au CNRS comme ailleurs, que l'issue de ma de-%
mande est douteuse, et que si un poste m'y était attribué, ce serait
aux dépens d'un chercheur plus jeune qui resterait sans poste.
Mais il est vrai aussi que cela libérerait mon poste à l'USTL au
bénéfice d'un autre. C'est pourquoi je n'ai pas de scrupule à
faire cette demande, et s'il le faut à revenir à la charge si elle
n'est pas acceptée cette année.

En tout état de cause, cette demande aura été pour moi
l'occasion d'écrire cette esquisse de programme, qui autrement sans
doute n'aurait jamais vu le jour. J'ai essayé d'être bref sans être
sibyllin et aussi, après coup, d'en faciliter la lecture et de
la rendre plus attrayante, en y adjoignant un sommaire. Si malgré
cela elle peut paraître longue pour la circonstance, je m'en excuse.
Elle me paraît courte pour son contenu, sachant bien que dix ans de
travail ne seraient pas de trop pour aller jusqu'au bout du moindre

\margin{128}

\clearpage
% Esquisse d'un programme (cote 119, copy 2): archive cote 119, site page 129 (PDF page 130), typescript page 51.
% typescript. Draft compared: esquisse/pages/esq_119_053.tex (copy 1, site 53).
% AI-transcribed by Claude Opus 5.5 (claude-opus-5-5, adaptive thinking, effort high), Message Batches API, batch esq119_full (msgbatch_015xqDHkTumDysRF1qngZhRJ), 2026-11-13; not yet checked by eye.
\scanpage{119/129}
\mspage{51}

des thèmes esquissés (à supposer qu'il y ait un ``bout''...), et
cent ans seraient peu pour le plus riche d'entre eux !

Derrière la disparité apparente des thèmes évoqués ici,
un lecteur attentif percevra comme moi une unité profonde. Celle-%
ci se manifeste notamment par une source d'inspiration commune,
la géométrie des surfaces, présente dans tous ces thèmes, au pre-%
mier plan le plus souvent. Cette source, par rapport à mon ``passé''
mathématique, représente un renouvellement, mais nullement une
rupture. Plutôt, elle montre le chemin d'une approche nouvelle
vers cette réalité encore mystérieuse, celle des ``\uline{motifs}'', qui me
fascinait plus que toute autre dans les dernières années de ce
passé (* ). Cette fascination ne s'est nullement évanouie,
elle fait partie plutôt de celle du plus brûlant pour moi de tous
les thèmes évoqués précédemment. Mais aujourd'hui je ne suis plus,
comme naguère, le prisonnier volontaire de tâches interminables,
qui si souvent m'avaient interdit de m'élancer dans l'inconnu,
mathématique ou non. Le temps des \uline{tâches} pour moi est révolu. Si
l'âge m'a apporté quelque chose, c'est d'être plus léger.

Janvier 1984

(* ) Voir à ce sujet mes commentaires dans ``l'Esquisse Thématique''
de 1972 jointe au présent rapport, dans la rubrique terminale
``divagations motiviques'', (loc. cit. pages 17-18).

\margin{129}

\clearpage
% Esquisse d'un programme (cote 119, copy 2): archive cote 119, site page 130 (PDF page 131), typescript page 52.
% typescript. Draft compared: esquisse/pages/esq_119_054.tex (copy 1, site 54).
% AI-transcribed by Claude Opus 5.5 (claude-opus-5-5, adaptive thinking, effort high), Message Batches API, batch esq119_full (msgbatch_015xqDHkTumDysRF1qngZhRJ), 2026-11-13; not yet checked by eye.
\scanpage{119/130}
\mspage{52}

(\textsuperscript{1}) L'expression ``hors de portée'' ici (et encore plus loin pour
une question toute différente), que j'ai laissé passer en allant
à l'encontre d'une réticence, me paraît décidément hâtive et sans
fondement. J'ai pu constater déjà en d'autres occasions que lors-%
que des augures (ici moi-même !) déclarent d'un air entendu
(ou dubitatif) que tel problème est ``hors de portée'', c'est là
au fond une affirmation entièrement subjective. Elle signifie
simplement, à part le fait que le problème est censé ne pas être
résolu encore, que celui qui parle est à court d'idées sur la
question, ou de façon plus précise sans doute, qu'il est devant
elle sans sentiment ni entrain, qu'elle ``ne lui fait rien'' et
qu'il n'a aucune envie de faire quelque chose avec elle - ce qui
souvent est une raison suffisante pour vouloir en décourager au-%
trui. Cela n'a pas empêché qu'à l'instar de M. de la Palisse, et
au moment même de succomber, les belles et regrettées conjectures
de Mordell, de Tate, de Chafarévitch étaient toujours réputées
``hors de portée'', les pauvres ! - D'ailleurs, dans les jours déjà
qui ont suivi la rédaction du présent rapport, qui m'a remis en
contact avec des questions dont je m'étais quelque peu éloigné au
cours de l'année écoulée, je me suis aperçu d'une nouvelle pro-%
priété remarquable de l'action extérieure d'un groupe de Galois
absolu sur le groupe fondamental d'une courbe algébrique, qui m'avait
échappée jusqu'à présent et qui sans doute constitue pour le moins
un nouveau pas en avant vers la formulation d'une caractérisation
algébrique de $\mathrm{Gal}(\bar{\underline{Q}}/\underline{Q})$. Celle-ci, avec la ``conjecture fondamen-%
tale'' (mentionnée au par. 3 ci-dessous) apparaît à présent comme
la principale question ouverte pour les fondements d'une ``géométrie
algébrique anabélienne'', laquelle depuis quelques années, repré-%
sente (et de loin) mon plus fort centre d'intérêt en mathématiques.

(\textsuperscript{2}) Je puis faire exception pourtant d'un autre ``fait'', du temps
où, vers l'âge de douze ans, j'étais interné au camp de concentra-%
tion de Rieucros (près de Mende). C'est là que j'ai appris, par une
détenue, Maria, qui me donnait des leçons particulières bénévoles,

\margin{130}

\clearpage
% Esquisse d'un programme (cote 119, copy 2): archive cote 119, site page 131 (PDF page 132), typescript page 53.
% typescript. Draft compared: esquisse/pages/esq_119_055.tex (copy 1, site 55).
% AI-transcribed by Claude Opus 5.5 (claude-opus-5-5, adaptive thinking, effort high), Message Batches API, batch esq119_full (msgbatch_015xqDHkTumDysRF1qngZhRJ), 2026-11-13; not yet checked by eye.
\scanpage{119/131}
\mspage{53}

la définition du cercle. Celle-ci m'avait impréssionné par sa sim-%
plicité et son évidence, alors que la propriété de ``rotondité
parfaite'' du cercle m'apparaissait auparavant comme une réalité mysté-%
rieuse au delà des mots. C'est à ce moment, je crois, que j'ai
entrevu pour la première fois (sans bien sûr me le formuler en
ces termes) la puissance créatrice d'une ``bonne'' définition ma-%
thématique, d'une \uline{formulation} qui décrit l'essence. Aujourd'hui
encore, il semble que la fascination qu'a exercé sur moi cette
puissance-là n'a rien perdu de sa force.

(\textsuperscript{3}) Plus généralement, au delà des variétés dites ``anabéliennes''
sur des corps de type fini, la géométrie algébrique anabélienne
\margin{/}(telle qu'elle s'est dégagé\hand{\unsure{e}} il y a quelques années) amène à une
description, en termes de groupes profinis uniquement, de la ca-%
tégorie des schémas de type fini sur la base absolue $\mathbb{Q}$ (voir
même $\mathbb{Z}$), et par là même, en principe, de la catégorie des sché-%
mas quelconques (par des passages à la limite convenables). Il
s'agit donc d'une construction ``qui fait semblant'' d'ignorer les
anneaux (tels que $\mathbb{Q}$, les algèbres de type fini sur $\mathbb{Q}$ etc) et
les équations algébriques qui servent traditionnellement à décri-%
re les schémas, en travaillant directement avec leurs topos éta-%
les, exprimables en termes de systèmes de groupes profinis. Un
grain de sel cependant : pour pouvoir espérer reconstituer un
schéma (de type fini sur $\mathbb{Q}$ disons) à partir de son topos étale,
qui est un invariant purement topologique, il convient de se pla-%
cer, non dans la catégorie des schémas (de type fini sur $\mathbb{Q}$ en
l'occurence), mais dans celle qui s'en déduit par ``localisation'',
en rendant inversibles les morphismes qui sont des ``homéomorphismes
universels'', i.e. qui sont finis, radiciels et surjectifs. Le
développement d'une telle traduction d'un ``monde géométrique''
(savoir celui des schémas, multiplicités schématiques etc) en
termes d'un ``monde algébrique'' (celui des groupes profinis, et
systèmes de groupes profinis décrivant des topos (dits ``étales'')
convenables) peut être considéré comme un aboutissement ultime

\margin{131}

\clearpage
% Esquisse d'un programme (cote 119, copy 2): archive cote 119, site page 132 (PDF page 133), typescript page 54.
% typescript. Draft compared: esquisse/pages/esq_119_056.tex (copy 1, site 56).
% AI-transcribed by Claude Opus 5.5 (claude-opus-5-5, adaptive thinking, effort high), Message Batches API, batch esq119_full (msgbatch_015xqDHkTumDysRF1qngZhRJ), 2026-11-13; not yet checked by eye.
\scanpage{119/132}
\mspage{54}

de la théorie de Galois, sans doute dans l'esprit même de Galois.
La sempiternelle question ``et pourquoi tout ça ?'' me paraît avoir
ni plus, ni moins de sens dans le cas de la géométrie algébrique
anabélienne en train de naître, que pour la théorie de Galois au
temps de Galois (ou même aujourd'hui, quand la question est posée
par un étudiant accablé ...), et de même pour le commentaire qui
va généralement avec : ``c'est bien général tout ça !''.

($^4$) On conçoit donc aisément qu'un groupe comme $S\ell(2,\mathbb{Z})$, avec
sa structure ``arithmétique'', soit une véritable machine à cons-%
truire des représentations ``motiviques'' de $\mathrm{Gal}(\bar{\mathbb{Q}}/\mathbb{Q})$ et de ses
sous-groupes ouverts, et qu'on obtient ainsi, au moins en prin-%
cipe, toutes les représentations motiviques qui sont de poids 1,
ou contenues dans un produit tensoriel de telles représentations
(ce qui en fait déjà un bon paquet !). J'avais commencé en 1981
à expérimenter avec cette machine dans quelques cas d'espèce,
obtenant diverses représentations remarquables de $\Gamma$ dans des
groupes $G(\mathbb{Z}^{\wedge})$, où $G$ est un schéma en groupes (pas nécessaire-%
ment réductif) sur $\mathbb{Z}$, en partant d'homomorphismes convenables
\[ S\ell(2,\mathbb{Z}) \longrightarrow G_o(\mathbb{Z}) , \]
où $G_o$ est un schéma en groupes sur $\mathbb{Z}$, et $G$ étant construit
à partir de là comme extension de $G_o$ par un schéma en groupes
convenable . Dans le cas ``tautologique'' $G_o = S\ell(2)_{\mathbb{Z}}$ , on trouve
pour $G$ une extension remarquable de $G\ell(2)_{\mathbb{Z}}$ par un tore de di-%
mension 2, avec une représentation motivique qui ``coiffe'' celles
associées aux corps de classes des extensions $\mathbb{Q}(i)$ et $\mathbb{Q}(j)$
(comme par hasard, les ``corps de multiplication complexe'' des deux
courbes elliptiques ``anharmoniques''). Il y a là un principe de
construction qui m'a semblé très général et très efficace, mais
je n'ai pas eu (ou pris) le loisir de le dévisser et le suivre
jusqu'au bout - c'est là un des nombreux ``points chauds'' dans le
programme de fondements de géométrie algébrique anabélienne (ou
de ``théorie de Galois'', version élargie) que je me propose de
développer. A l'heure actuelle, et dans un ordre de priorité sans

\margin{132}

\clearpage
% Esquisse d'un programme (cote 119, copy 2): archive cote 119, site page 133 (PDF page 134), typescript page 55.
% typescript. Draft compared: esquisse/pages/esq_119_057.tex (copy 1, site 57).
% AI-transcribed by Claude Opus 5.5 (claude-opus-5-5, adaptive thinking, effort high), Message Batches API, batch esq119_full (msgbatch_015xqDHkTumDysRF1qngZhRJ), 2026-11-13; not yet checked by eye.
\scanpage{119/133}
\mspage{55}

doute très provisoire, ces points sont :

a) Construction combinatoire de la Tour de Teichmüller.

b) Description du groupe des automorphismes de la compactifi-%
cation profinie de cette tour, et réflexion sur une caractérisa-%
tion de $\Gamma = \mathrm{Gal}(\bar{\mathbb{Q}}/\mathbb{Q})$ comme sous-groupe de ce dernier.

c) La ``machine à motifs'' $S\ell(2,\mathbb{Z})$ et ses variantes.

d) Le dictionnaire anabélien, et la conjecture fondamentale
(qui n'est peut-être pas si ``hors de portée'' que ça !). Parmi
les points cruciaux de ce dictionnaire, je prévois le ``paradigme
profini'' pour les corps $\mathbb{Q}$ (cf. b)), $\mathbb{R}$ et $\mathbb{C}$, dont une formu-%
lation plausible reste à dégager, ainsi qu'une description des sous-%
groupes d'inertie de $\Gamma$ , par où s'amorce le passage de la caracté-%
ristique zéro à la caractéristique $p > 0$, et à l'anneau absolu $\mathbb{Z}$.

e) Problème de Fermat.

($^5$) Je signalerai cependant un travail plus délicat (mis à part
le travail signalé en passant sur les complexes cubiques), sur
l'interprétation combinatoire des cartes régulières associées aux
sous-groupes de congruence de $S\ell(2,\mathbb{Z})$. Ce travail a été développé
surtout en vue d'exprimer l'opération ``arithmétique'' de $\Gamma = \mathrm{Gal}(\bar{\mathbb{Q}}/\mathbb{Q})$
sur ces ``\uline{cartes de congruence}'', laquelle se fait, essentiellement,
par l'intermédiaire du caractère cyclotomique de $\Gamma$. Un point de
départ a été la théorie combinatoire du ``bi-icosaèdre'' développée
dans un cours C4 à partir de motivations purement géométriques,
et qui (il s'est avéré par la suite) permet d'exprimer commodément
l'opération de $\Gamma$ sur la catégorie des cartes icosaédrales (i.e.
les cartes de congruence d'indice 5).

($^6$) Signalons à ce propos que les classes d'isomorphie d'espaces
modérés compacts sont les mêmes que dans la théorie ``linéaire par
morceaux'' (qui n'est \uline{pas}, je le rappelle, une théorie modérée).
C'est là, en un sens, une réhabilitation de la ``Hauptvermutung'',
qui n'est ``fausse'' que parce que, pour des raisons historiques
qu'il serait sans doute intéressant de cerner de plus près, les

\margin{133}

\clearpage
% Esquisse d'un programme (cote 119, copy 2): archive cote 119, site page 134 (PDF page 135), typescript page 56.
% typescript. Draft compared: esquisse/pages/esq_119_058.tex (copy 1, site 58).
% AI-transcribed by Claude Opus 5.5 (claude-opus-5-5, adaptive thinking, effort high), Message Batches API, batch pilotnext (msgbatch_01A9sdu8SYq23bEtqX9tUncQ), 2026-11-13; not yet checked by eye.
\scanpage{119/134}
\mspage{56}

fondements de topologie utilisés pour la formuler n'excluaient
pas les phénomènes de sauvagerie. Il va (je l'espère) sans dire
que la nécessité de développer de nouveaux fondements pour la
topologie ``géométrique'' n'exclut nullement que les phénomènes en
question, comme toute chose sous le ciel, ont leur raison d'être
et leur propre beauté. Des fondements plus adéquats ne suppri-%
meront pas ces phénomènes, mais nous permettront de les situer à
leur juste place, comme des ``cas limites'' de phénomènes de ``vraie''
topologie..

($^7$) En fait, pour reconstituer ce système d'espaces
\[ (i_o,\ldots,i_n) \longmapsto X_{i_o,\ldots,i_n} \]
contravariant sur Drap(I) (pour l'inclusion des drapeaux), il
suffit de connaître les $X_i$ (ou ``\uline{strates déployées}'') et les $X_{i,j}$
(ou ``\uline{tubes de raccord}'') pour $i,j \in I$, $i < j$, et les morphismes
$X_{i,j} \longrightarrow X_i$ (qui sont des inclusions ``bordantes'') et $X_{i,j} \rightarrow X_i$
(qui sont des fibrations propres, dont les fibres $F_{i,j}$ sont
appelées ``\uline{fibres de raccord}'' pour les strates d'indice $i$ et $j$).
Dans le cas d'une multiplicité modérée cependant, il faut connaî-%
tre de plus les ``\uline{espaces de jonction}'' $X_{i,j,k}$ $(i < j < k)$ et ses
morphismes dans $X_{i,j}$ , $X_{j,k}$ , et surtout $X_{i,k}$ , s'insérant dans
le diagramme commutatif hexagonal suivant, où les deux carrés de
droite sont cartésiens, les flèches $\hookrightarrow$ sont des immersions (pas
nécessairement des plongements ici), et les autres flèches sont
des fibrations propres :

\begin{diagramblock}
\begin{tikzcd}
 & X_{ik} \arrow[rr, hook] \arrow[dl] & & X_k & \\
X_i & & X_{ijk} \arrow[ul, hook] \arrow[rr, hook] \arrow[dl] & & X_{jk} \arrow[ul, hook] \arrow[dl] \\
 & X_{ij} \arrow[ul] \arrow[rr, hook] & & X_j &
\end{tikzcd}
\end{diagramblock}

(NB. Ce diagramme définit $X_{ijk}$ en termes de $X_{ij}$ et $X_{jk}$ sur $X_j$,
mais non la flèche $X_{ijk} \longrightarrow X_{ik}$ , car $X_{ik} \longrightarrow X_k$ n'est \uline{pas} né-%
cessairement un plongement.)

\margin{134}

\clearpage
% Esquisse d'un programme (cote 119, copy 2): archive cote 119, site page 135 (PDF page 136), typescript page 57.
% typescript. Draft compared: esquisse/pages/esq_119_059.tex (copy 1, site 59).
% AI-transcribed by Claude Opus 5.5 (claude-opus-5-5, adaptive thinking, effort high), Message Batches API, batch esq119_full (msgbatch_015xqDHkTumDysRF1qngZhRJ), 2026-11-13; not yet checked by eye.
\scanpage{119/135}
\mspage{57}

Dans le cas des espaces modérés stratifiés proprement dits (qui
ne sont pas des multiplicités à proprement parler) on peut expri-%
mer de façon commode le ``déployement'' de cette structure, i.e.
le système des espaces $X_{i_o,\ldots,i_n}$ , en termes de l'espaces mo-%
déré $X_*$ somme des $X_i$ , qui est munie d'une \uline{structure d'objet}
\uline{ordonné} (dans la catégorie des espaces modérés) ayant comme gra-%
phe $X_{**}$ de la relation d'ordre la somme des $X_{i,j}$ et des $X_i$
(ces derniers constituant la diagonale). Parmi les propriétés
essentielles de cette structure ordonnée, relevons seulement
ici que $pr_1 : X_{**} \longrightarrow X_*$ est une fibration (localement trivia-%
le) propre, et $pr_2 : X_{**} \longrightarrow X_*$ est une plongement ``bordant''.
On a une interprétation analogue du déployement d'une multipli-%
cité modérée stratifiée, en termes d'une \uline{structure de catégorie}
(remplaçant une simple structure ordonnée) ``au sens multipli-%
cités modérées'', dont l'application de composition est donnée
par les morphismes $X_{ijk} \longrightarrow X_{ik}$ ci-dessus.

\margin{135}

\end{document}
